\documentclass[11pt,oneside,openany,a4paper]{report}
\usepackage[T1]{fontenc}
\usepackage[utf8]{inputenc}
\usepackage{lmodern,microtype}
\usepackage[a4paper,left=26mm,right=26mm,top=25mm,bottom=25mm,headheight=15pt,headsep=7mm]{geometry}
\usepackage{amsmath,amssymb,amsthm,mathtools,bm,mathrsfs}
\usepackage{booktabs,array,tabularx,longtable,ragged2e}
\usepackage{enumitem,needspace,fancyhdr,xcolor,titlesec,graphicx}
\usepackage{hyperref,xurl,bookmark}
\definecolor{inkblue}{RGB}{28,54,78}
\hypersetup{colorlinks=true,linkcolor=inkblue,citecolor=inkblue,urlcolor=inkblue,pdfauthor={},bookmarksnumbered=true}
\pagestyle{fancy}\fancyhf{}
\fancyhead[L]{\small HSS approximation: research record}
\fancyhead[R]{\small\thepage}
\renewcommand{\headrulewidth}{.3pt}
\fancypagestyle{plain}{\fancyhf{}\fancyfoot[C]{\thepage}\renewcommand{\headrulewidth}{0pt}}
\titleformat{\chapter}[display]{\normalfont\huge\bfseries\color{inkblue}\raggedright}
 {\large\normalfont\scshape Chapter \thechapter}{8pt}{\Huge}
\titlespacing*{\chapter}{0pt}{0pt}{19pt}
\titleformat{\section}{\normalfont\Large\bfseries\color{inkblue}}{\thesection}{.7em}{}
\titleformat{\subsection}{\normalfont\large\bfseries}{\thesubsection}{.7em}{}
\setlength{\parindent}{0pt}
\setlength{\parskip}{4pt plus 1pt minus .5pt}
\setlength{\emergencystretch}{3em}
\setlist{topsep=4pt,itemsep=2pt,parsep=0pt}
\allowdisplaybreaks[2]
\clubpenalty=10000\widowpenalty=10000\displaywidowpenalty=10000
\setcounter{tocdepth}{1}\setcounter{secnumdepth}{2}
\numberwithin{equation}{chapter}
\newtheorem{theorem}{Theorem}[chapter]
\newtheorem{lemma}[theorem]{Lemma}
\newtheorem{proposition}[theorem]{Proposition}
\newtheorem{corollary}[theorem]{Corollary}
\theoremstyle{definition}
\newtheorem{definition}[theorem]{Definition}
\newtheorem{assumption}[theorem]{Assumption}
\newtheorem{conjecture}[theorem]{Conjecture}
\newtheorem{obstruction}[theorem]{Obstruction}
\theoremstyle{remark}
\newtheorem{remark}[theorem]{Remark}
\newtheorem{example}[theorem]{Example}
\newtheorem{auditnote}[theorem]{Qualification}
\newtheorem{checkpoint}[theorem]{Open checkpoint}
\newcommand{\sourcefile}[1]{{\small\nolinkurl{#1}}}
\newcommand{\R}{\mathbb R}\newcommand{\E}{\mathbb E}\newcommand{\PP}{\mathbb P}\newcommand{\Prb}{\mathbb P}
\newcommand{\T}{\mathsf T}\newcommand{\Id}{I}
\newcommand{\norm}[1]{\left\lVert#1\right\rVert}
\newcommand{\F}[1]{\left\lVert#1\right\rVert_F}
\newcommand{\N}[1]{\left\lVert#1\right\rVert}
\newcommand{\op}[1]{\left\lVert#1\right\rVert_{\mathrm{op}}}
\newcommand{\ip}[2]{\left\langle#1,#2\right\rangle_F}
\DeclareMathOperator{\tr}{tr}\DeclareMathOperator{\rank}{rank}
\DeclareMathOperator{\diag}{diag}\DeclareMathOperator{\range}{range}
\DeclareMathOperator{\ran}{range}\DeclareMathOperator{\OPT}{OPT}
\DeclareMathOperator{\HSS}{HSS}\DeclareMathOperator{\dist}{dist}
\DeclareMathOperator{\Var}{Var}\DeclareMathOperator{\supp}{supp}
\DeclareMathOperator{\orth}{orth}\DeclareMathOperator*{\argmin}{arg\,min}
\DeclareMathOperator{\spanop}{span}\DeclareMathOperator{\Cov}{Cov}
\DeclareMathOperator{\Proj}{Proj}\DeclareMathOperator{\vecop}{vec}
\newcommand{\HH}{\mathcal H}\newcommand{\SSS}{\mathcal S}
\newcommand{\UC}{\mathcal U}\newcommand{\VC}{\mathcal V}\newcommand{\err}{\mathcal E}
\newcolumntype{Y}{>{\RaggedRight\arraybackslash}X}
\pdfstringdefDisableCommands{\def\Psi{Psi}\def\Omega{Omega}\def\epsilon{epsilon}\def\varepsilon{epsilon}\def\sigma{sigma}\def\tau{tau}\def\alpha{alpha}\def\beta{beta}\def\mathcal#1{#1}\def\mathbb#1{#1}\def\operatorname#1{#1}\def\norm#1{norm(#1)}\def\sourcefile#1{#1}}

\hypersetup{pdftitle={Obstructions in Shared-Probe HSS Approximation}}
\begin{document}
\hypersetup{pageanchor=false}
\begin{titlepage}
\vspace*{10mm}
{\color{inkblue}\large\scshape Technical manuscript\par}
\vspace{17mm}
{\color{inkblue}\Huge\bfseries Obstructions in\par Shared-Probe HSS\par Approximation\par}
\vspace{9mm}
{\Large\raggedright Counterexamples, failed estimates, repaired gaps,\\ and the remaining barriers\par}
\vspace{13mm}
{\large Research snapshot: 20 September 2026\par}
\vspace{15mm}
\hrule\vspace{5mm}
\textbf{Status.} The general conjecture remains open in this research record. This manuscript records what failed, why it failed, and the exact scope of each obstruction.
\vspace{5mm}\hrule
\vfill
Prepared from the supplied proof manuscript and the subsequent proof,
audit, and numerical records.\\[4pt]
Real matrices. Fixed inputs. Exact arithmetic in theorem statements.\\
Squared Frobenius error and explicit conditioning conventions.
\end{titlepage}
\pagenumbering{roman}\hypersetup{pageanchor=true}
\chapter*{Abstract}
\addcontentsline{toc}{chapter}{Abstract}
This manuscript consolidates the obstructions encountered in an attempt to obtain an $O(k)$ matrix--vector query algorithm for quasi-optimal HSS approximation. The desired guarantee is an expected squared Frobenius error bounded by $C L\OPT_k(A)$, with final HSS rank at most $k$ and polynomial additional computation. Neither a proof nor a disproof of that conjecture has been obtained.

The record distinguishes failures of proof estimates from failures of specified algorithms. It includes adaptive Gaussian conditioning, inverse-moment growth on exact inputs, noncommutative and matrix-embedding obstructions in refitting, actual noisy-input counterexamples for earlier basis learners, failure of screened residual estimates after dropping the final projection, and failures of global interpolation selectors even with exact leaf information. It also records why local independent-source theorems, independent pilots, and finite numerical tests do not yet imply the general shared-probe result. Historical refitting obstructions are presented with their resolutions: the independent fixed-space refit now has a depth-uniform factor $1607$.

Later counterexamples distinguish the old full-range learner, separately trained two-view hierarchies, and the current synchronized learner. Negative results for the first two do not refute the third. The concluding status map identifies the remaining adaptive signed-error estimate and the separate efficient rank-reduction requirement. New observations that had not completed independent review when research stopped are explicitly segregated from established results.

\chapter*{Scope and conventions}
\addcontentsline{toc}{chapter}{Scope and conventions}
This is a technical synthesis of \sourcefile{hss_complete_proof_chain.tex} and the subsequent local proof and audit notes, frozen on 20 September 2026 at the user's instruction to stop research. It is not a claim that every retained argument has received external peer review. Here ``reviewed'' or ``audited'' means that a corresponding internal mathematical audit was recorded. The supplied original manuscript has not been overwritten.

The motivating public benchmark is Amsel et al.\ \cite{amsel2025}: expected squared Frobenius error $O(L)\OPT_k(A)$ using $O(kL)$ queries. The open improvement seeks to remove the depth factor from the query count, while retaining the approximation guarantee, output rank, and computational requirement. Levitt--Martinsson's reused-sketch construction is relevant background, but its discussion of coarse-level sketch quality does not provide the missing all-input approximation theorem \cite{levitt2024}. This manuscript makes no exhaustive literature-absence claim.

All probabilistic assertions concern a fixed input chosen before the random probes unless explicitly stated otherwise. A data-dependent competing interpolant is permitted in an argument about the failure of an estimator; it is not silently reinterpreted as a fixed adversarial input. Positive-probability event constructions are distinguished from lower bounds whose probabilities remain uniform under a parameter limit.

We use $\OPT_k(A)=\inf_{B\in\mathcal H_k}\|A-B\|_F^2$ and write $L$ for tree depth, counting the nonroot cut levels under the local convention. Physical leaf sizes differ between candidate algorithms and are specified with each result. Node projectors onto coordinate sets and projectors onto retained learned spaces are different operators; they are disambiguated locally.

The companion manuscript, \emph{The Current Proof Route for Shared-Probe HSS Approximation}, specifies the surviving synchronized construction and its exact outstanding estimate. The present document groups repeated attempts by mathematical obstruction rather than reproducing a chronological conversation.

\paragraph{Status vocabulary.}
A \emph{counterexample} refutes the particular quantified claim stated beside it. A \emph{gap} is a premise still needed by an otherwise valid implication. A \emph{repaired obstruction} explains why an earlier argument failed and what later theorem replaces it. A \emph{numerical diagnostic} is finite evidence under its reported sampling design. A \emph{working observation} was not promoted to an audited theorem before the stop point.

\begingroup\small
\tableofcontents
\endgroup
\clearpage
\pagenumbering{arabic}
\chapter{Original proof-chain obstructions}
% Body contribution: original-manuscript obstruction taxonomy.
% Source labels below are printed identifiers from the original manuscript,
% not cross-document references. No preamble or document environment.

\section{Obstructions in the original manuscript: a logical taxonomy}
\label{origobs:taxonomy}

This chapter consolidates the negative results and proof limitations in
\sourcefile{hss_complete_proof_chain.tex}.  Its purpose is to specify what
each obstruction actually excludes.  A counterexample to a proposed
potential, a failure of a particular learned space, and an information
theoretic lower bound are different statements.  The original manuscript
contains examples of the first two kinds, but no impossibility theorem
for the general $O(k)$-query HSS approximation problem.

Throughout, $A$ is fixed before the probes, $\mathcal H_k$ is the HSS
rank-$k$ class on the prescribed finite coordinate tree, and
\[
 e_k(A)=\inf_{B\in\mathcal H_k}\|A-B\|_F^2.
\]
For a learned frozen coefficient space $\mathcal S$, write
$d_{\mathcal S}(A)^2=\|(I-P_{\mathcal S})A\|_F^2$.  When
$\mathcal S\subseteq\mathcal H_r$, a theorem relative to
$d_{\mathcal S}(A)$ is a statement about that particular learned space,
not automatically a theorem relative to $e_k(A)$.  In particular,
$d_{\mathcal S}(A)^2\ge e_k(A)$ when $\mathcal S\subseteq\mathcal H_k$;
the inequality points in the wrong direction for replacing the
fixed-space comparator in an upper bound.

\paragraph{Status convention.}
We distinguish \emph{counterexamples}, which disprove a stated universal
claim; \emph{missing premises}, which identify a necessary input to an
otherwise valid implication; and \emph{repaired historical gaps}, whose
resolution appears later in the same original manuscript.  The appendix
labelled \texttt{base:new:route-history} records 66 routes, and
\texttt{base:new:obstruction-ledger} records their logical scopes.
Some of its final paragraphs preserve an earlier, unresolved status for
fixed-space refitting.  That historical status is superseded by Chapters
\emph{Frozen-stream feedback}, \emph{Coherent drift}, and
\emph{Uniform and near-optimal fixed-space refitting}
(\texttt{ch:feedback}, \texttt{ch:drift}, \texttt{ch:uniform}).
The established endpoint is the factor $1607$ with $8k$ independent
refit queries.  The basis-learning premise remains separate.

\subsection{Comparators, query accounting, and exactness are separate interfaces}
\label{origobs:interfaces}

Fresh queries at every level supply a legitimate comparison algorithm,
but an $O(kL)$-query proof does not become an $O(k)$-query reuse proof by
omitting the freshness hypothesis.  Likewise an $O(L)$ learning factor
and an $O(L)$ coefficient-refitting factor compose to $O(L^2)$, not to
the desired $O(L)$.  These accounting restrictions are explicit in the
original query ledger, \texttt{base:new:ledger}.

Exact recovery on fixed HSS inputs is also weaker than residual-relative
stability.  An exact trajectory may be discontinuous across a
rank-changing stratum, and its inverse constants may be unbounded near
that stratum.  Canonical completion removes arbitrary choices of zero
singular vectors; it does not provide a quantitative perturbation
theorem.  The original raw learner's exact-input result and its first
two compression estimates therefore do not supply an all-depth
approximation induction (\texttt{base:sec:exact-recovery},
\texttt{base:sec:finite}, \texttt{base:sec:second}).

For a known bicriteria output $C$, exact offline HSS rank reduction is a
valid query-only operation: if $\widehat B$ minimizes
$\|C-B\|_F$ over $\mathcal H_k$, comparison with a best $B_*$ for $A$
gives
\[
 \|A-\widehat B\|_F\le2\|A-C\|_F+\sqrt{e_k(A)}.
\]
This does not assert polynomial-time computability or authorize replacing
the exact minimizer by a practical recompression routine with a different
factor.  Nor does a lower bound for the unprocessed learned-space
candidate automatically persist after arbitrary nonlinear postprocessing.
These qualifications are part of \texttt{ch:bicriteria} and
\texttt{ch:synthesis}.

\section{Gaussian beginnings do not give Gaussian adaptive regressions}
\label{origobs:adaptive-inverses}

\subsection{What ambient Gaussian compression really does prove}

Let $G\in\mathbb R^{d\times s}$ be standard Gaussian, $s>d+1$, and let
$W\in\mathbb R^{d\times r}$ be any measurable isometry, even one chosen
from $G$.  The valid pathwise comparison is
\[
 (W^TGG^TW)^{-1}\preceq W^T(GG^T)^{-1}W,
 \qquad
 \mathbb E\|(W^TG)^\dagger\|_F^2\le\frac{d}{s-d-1}.
\]
The inverse is controlled by the \emph{ambient} dimension $d$, not merely
the retained dimension $r$.  The proof is a Schur-complement comparison
followed by the Gaussian inverse trace formula.  Thus adaptivity alone
does not invalidate every inverse estimate; it invalidates silently
replacing an adaptively compressed design by a fresh $r$-row Gaussian
matrix.  The source is Lemma \texttt{base:lem:inverse-compression}.

\subsection{Nested full rank can coexist with an infinite inverse moment}

Proposition \texttt{base:prop:nested-obstruction} constructs a measurable
binary hierarchy of rank-$k$ orthogonal compressions of a square standard
Gaussian $G\in\mathbb R^{8k\times8k}$.  Its cumulative isometry
$W(G)\in\mathbb R^{8k\times2k}$ has full-row-rank compressed design
$W^TG$ almost surely, but
\[
 \mathbb E\|(W^TG)^\dagger\|_F^2=\infty.
\]
Choose a left singular vector $u$ for the smallest singular value of $G$.
At each binary stage retain the restrictions of $u$ in the appropriate
children, padding orthonormally as needed.  The resulting nested range
contains $u$, so compression retains the smallest eigenvalue of $GG^T$.
The reciprocal squared distance of the last Gaussian row from the span
of its predecessors has the law $Z^{-2}$ with $Z\sim N(0,1)$ and infinite
mean.  It lower-bounds the relevant inverse norm.

The input to this construction is the probe itself.  It is not claimed
to be the SVD hierarchy reached by the actual learner on a fixed $A$.
Its conclusion is exact: initial Gaussianity, orthonormality and nesting
alone are insufficient hypotheses.  With more physical rows than probe
columns, unrestricted adaptive nesting can even preserve a vector in
the exact left kernel.

The fractional-moment version is stronger.  In Proposition
\texttt{base:aug:prop:all-fractional}, let $G\in\mathbb R^{d\times s}$
with $d>s$, choose unit vectors $u_0\in\ker G^T$ and $u_1$ in a largest
left singular direction, and put
\[
 R=\|G\|_F,\qquad \varepsilon=\exp(-R^4),\qquad
 w=\sqrt{1-\varepsilon^2}\,u_0+\varepsilon u_1.
\]
Then $w^TG$ is nonzero almost surely, while for every $p>0$
\[
 \|(w^TG)^\dagger\|^{2p}
 \ge \exp(2pR^4)R^{-2p},\qquad
 \mathbb E\|(w^TG)^\dagger\|^{2p}=\infty.
\]
The chi-density of $R$ cannot offset the quartic exponential.  This
excludes a universal fractional inverse theorem for arbitrary adaptive
compressions; it does not disprove a weak-tail theorem for the actual
HSS rule.

\subsection{Separate energy and inverse moments do not price their product}

Example \texttt{base:ex:correlated-inverse} uses a standard Gaussian row
$g$ and $p_\varepsilon=\Pr\{\|g\|\le\varepsilon\}$.  Set
\[
 f(g)=p_\varepsilon^{-1/2}\mathbf1_{\{\|g\|\le\varepsilon\}}
       \frac{g}{\|g\|}.
\]
Although $\mathbb E\|f(g)\|^2=1$, the regression cost satisfies
\[
 \mathbb E|f(g)g^\dagger|^2
 =\mathbb E[\|g\|^{-2}\mid\|g\|\le\varepsilon]
 \ge\varepsilon^{-2}.
\]
Even a finite marginal inverse moment does not justify factorizing this
dependent product.  The same issue survives replacing a conditional
Wishart covariance by its expectation: inverse moments and nonlinear
regression risks do not commute with that averaging.  The correct
Wishart interface in \texttt{base:new:wishart} first identifies the
original conditional Gaussian law; it is not a freshness argument.

\subsection{Stein corrections, orthogonal mixing, and ridge do not give a shortcut}

The regularized integration-by-parts calculation in
\texttt{base:sec:closure} separates a fresh regression cost from an
adaptation derivative.  Its sign is not universally favorable.  If a
projector selects the $r$ weakest left singular directions of an
$m$-row design, it selects the $r$ largest inverse eigenvalues.  Their
sum is strictly larger than $r/m$ times the full inverse trace, so the
risk exceeds the fresh coefficient.  This is a counterexample to a
sign inference, not a description of the actual HSS selector.  The
zero-ridge limit additionally requires justified integrability; a
regularized identity alone does not provide it.

Common right orthogonal mixing of a data/design pair,
$(Y,\Omega)\mapsto(YO,\Omega O)$, leaves the regression and innovation
Gram unchanged.  It merely changes query coordinates and does not
refresh the estimator.  Positive ridge in the discrepancy has a
different defect: it generally biases even exact inputs such as the
identity, whose nullified innovations vanish.  Exact recovery is a
necessary requirement for a residual-relative theorem when
$e_k(A)=0$.  These are the two modifications discussed in
\texttt{base:sec:additional-obstructions}; neither settles the reuse
problem.

\section{Validation and global fitting: exact implications with missing premises}
\label{origobs:validation}

The validation route in \texttt{base:sec:validation} is a conditional
amplification theorem.  It can convert a suitable weak tail for the
\emph{actual} candidate error into an expected-error guarantee by making
independent candidates and testing them with independent probes.  It
does not establish that weak tail.  The distinction is visible even in
one dimension: if $F=Z^{-2}$ for a standard normal $Z$, then
$\Pr(F>t)\le\sqrt{2/\pi}\,t^{-1/2}$; the minimum of three independent
copies has finite expectation, whereas the minimum of two still does
not.  Conversely, a statement saying only that a candidate succeeds
with probability $0.9$ is insufficient.  A loss equal to $M$ with
probability $0.9$ and $TM$ otherwise has
\[
 \mathbb E\min_{1\le i\le a}F_i
 =M\bigl[1+(T-1)10^{-a}\bigr],
\]
which is unbounded in $T$ for every fixed number $a$ of candidates.
Clipping or a fixed-probability success statement cannot substitute for
tail control relative to the unknown optimum.

Training consistency has a separate uniformity defect.  If both
training widths are below the ambient dimension, choose
$u\in\ker\Psi^T$ and $v\in\ker\Omega^T$.  The rank-one matrix $uv^T$ is
invisible to both sketches.  This rules out a residual certificate
uniform over all sketch-dependent candidates, even within HSS rank one.
Because this invisible matrix is chosen after the probes, the example
is not a fixed-input information lower bound.

\subsection{Identifiability does not provide uniform nonlinear conditioning}

The positive identifiability result in
\texttt{base:con:identifiability} must be retained alongside its
limitations.  For a fixed HSS matrix whose prescribed cut ranks are
saturated at $k$, sufficiently wide two-sided Gaussian observations
(width at least $3k$ in the stated theorem) identify its cut spaces and
cores almost surely.  The competing HSS matrices may depend on the
observations.  A local smooth inverse exists on the corresponding
full-rank, spectral-gap charts.  Its radius and Lipschitz constant are
random, input-dependent quantities; the argument does not integrate
them to obtain uniform expected approximation.

Nonsaturated inputs behave differently.  If a unit vector $u$ lies in
the common nullspace of $\Omega^T$ and $\Psi^T$, then $I$ and
$I-uu^T$ have identical observations.  Their off-diagonal cut ranks are
at most one, but $I-uu^T$ has smaller Frobenius norm.  Thus a
minimum-norm completion can fail on the exact identity input.  The
manuscript's lexicographic cut-rank selector addresses exact rank
selection, not noisy stability.

The finite-family route in \texttt{ch:finite} has a different
obstruction.  Even leaf-block-diagonal members contain a Euclidean
family of dimension proportional to $Nk$.  A generic net therefore
has logarithmic cardinality of that order, and generic finite-model
selection brings an ambient-size penalty.  This is a limitation of
the net argument: those same diagonal blocks can be recovered with
$2k$ structured forward products.  Likewise, direct global least
squares can be injective without a proved depth-uniform lower bound
on the measurement map over the entire coefficient space.

\subsection{A concrete failure of depth-weighted nuclear selection}

Theorem \texttt{nuclear:thm:failure} is an actual selector
counterexample, not merely a conditioning warning.  Divide the root
into two halves $I,J$ of size $n$, let the two probe widths be $q$,
and fix
$B[I,J]=uv^T$ with constant unit vectors $u,v$; all other blocks vanish.
Let $P$ project onto $\operatorname{range}\Psi_I$ and $Q$ onto
$\operatorname{range}\Omega_J$.  The sketch-dependent competitor
\[
 C=Pu\,v^T+u\,v^TQ-Pu\,v^TQ
   =(Pu)v^T+(u-Pu)(Qv)^T
\]
has rank at most two and the same two observations as $B$.  Setting
$a=\|Pu\|$ and $b=\|Qv\|$ gives
\[
 \|C\|_F^2=a^2+b^2-a^2b^2,\qquad
 \mathbb E(a^2+b^2)=2q/n.
\]
Hence $\|C\|_F^2<1/2$ on an event of probability at least $1-4q/n$.

For any nonnegative depth weights, not all zero, the sum of the
nuclear norms of both orientations of all cuts at a given depth is
$2\sqrt m$ for $B$, where $m$ is the number of nodes at that depth.
For $C$ it is at most $2\sqrt{2m}\|C\|_F$.  Thus the weighted tree
nuclear penalty is strictly smaller for $C$ on that event.  For
$k\ge2$ and $n>4q$, either exact completion by this penalty or
least-squares fitting plus a positive multiple of it fails to return
the fixed exact input $B$.  The conclusion excludes this specified
regularizer, with its stated depth weights; it does not exclude all
global fitting objectives.

\section{Signed quotient energy and the failure of scalar bookkeeping}
\label{origobs:quotient}

In the original refit analysis, the innovation and mismatch state
removes a common coefficient that has no effect on the next basis
selection.  Its energy may be written
\[
 V=\|E_Y\|_F^2+\|E_Z\|_F^2+\|\Delta\|_F^2.
\]
Compression dissipates this energy exactly:
\[
 V'=V-\|U_\perp^TE_Y\|_F^2-\|V_\perp^TE_Z\|_F^2
        -\|\Delta_{22}\|_F^2.
\]
At a merge, however, the only surviving interaction is signed:
\[
 V_{\rm parent}=V_1+V_2-2\langle F_Y,F_Z^T\rangle.
\]
The scalar-weight experiment in \texttt{base:new:energy} is conclusive
for its proposed class: compression requires a weight at least one,
whereas the state $F_Z^T=-F_Y$ makes the merge require a weight at most
one half.  No single scalar weight contracts both maps on every
admissible state.

Even a successfully contracting quotient is not an output theorem.
The pair $Y=C\Omega$, $Z=C^T\Psi$ has zero quotient but can carry a
nonzero common coefficient $C$.  Emitted coefficients and the terminal
root must be recorded separately.  Similarly, inactive error is a
reconstruction charge, not future active forcing; charging it at every
later level repeatedly counts the same physical loss
(\texttt{base:sec:global}, \texttt{base:new:ledger-interface}).

The reachable rank-two witness in
\texttt{base:new:scalar-storage} shows that adding a scalar residual
storage does not cure the problem.  With $P=e_1e_1^T$,
$Q=I-P$, use the completion
\[
 \begin{pmatrix}P&-Q\\Q&P\end{pmatrix}.
\]
At $k=2,q=8$, the source state can have $Y=F=C=D=0$ and
$R=J_0=e_1e_1^T$.  A fresh row $\xi$ then produces
$C'=\xi_2e_1e_2^T$, $F'=-\xi_2e_2g_1$, and
$R'=(1-\xi_2^2)e_1e_1^T$, with
$\mathbb E\xi_2^2=1/4$.  The proposed scalar output-plus-storage
inequality fails for every scalar weight, although the actual output
$Y'$ is zero.  Directional matrix storage resolves this particular
witness; it does not turn its restricted geometric hypotheses into a
theorem for arbitrary learned bases.

\subsection{Independent one-pass estimates are not feedback estimates}

A centered one-pass map $K_G$ can satisfy
$\mathbb E\|K_Gb\|^2\le\rho\|b\|^2$ for deterministic $b$, or for
$b$ independent of $G$.  In the coupled equations
$a=a_0-K_Gb$, $b=b_0-K_Ha$, the solved inputs are functions of both
streams.  Iterating the independent-input estimate is therefore an
unjustified substitution.

Chapter \texttt{ch:halfsolve} supplies a finite-dimensional
counterexample even with independent half solvers.  Let $S$ be the
nilpotent shift on $\mathbb R^h$ and
$P=\operatorname{diag}(1,\ldots,1,0)$.  Independently choose
$\xi,\eta$ equal to $\pm1$ with respective total probabilities $a,b$,
and zero otherwise; set $C=\xi S$, $D=\eta P$.  Both maps are centered
contractions with the appropriate one-pass second-moment bounds, and
$CD=\xi\eta S$ is nilpotent.  Nevertheless
\[
 \mathbb E\|(I-CD)^{-1}e_1\|^2=1+ab(h-1).
\]
The same rare signs activate every power of the feedback loop; their
probabilities do not multiply afresh at each traversal.  A residual
injection gives coupled physical error
$a[1+ab(h-1)]$ while preserving the stated individual half-solve
guarantees.  This abstract obstruction does not prove failure of the
specific HSS refit or of a positive global least-squares decoder.

\section{Matrix Carleson control does not remove anticipative alignment}
\label{origobs:embedding}

The restart theorem in \texttt{base:new:carleson} provides an adapted
conditional future-response estimate.  Its valid proof retains the
full cross covariance tensor and deterministic mean transport
(\texttt{base:con:restart-cone}).  That is a repaired historical gap:
vanishing of a scalar cross trace alone was insufficient to preserve
the required covariance cone.  The source decomposition also had to
include discarded directions \emph{inside} a nonleaf sibling.
Example \texttt{base:con:interior-example} invalidates their omission.
The repaired coarse off-chain budget remains valid, with the distinct
$q-2k$ and $q-3k$ denominators kept in their respective modules.

These repairs do not imply a dimension-free embedding for an
anticipatively selected terminal source.  The distinction between an
adapted restart and a source using future probe information is
essential.  The original embedding route yields logarithmic
rank/time factors; an isotropy assertion alone cannot delete them.

\subsection{The Haar--harmonic obstruction}

The construction in \texttt{base:new:abstract-obstruction} takes
$m\ge4$, ambient dimension at least $2m$, and successive Haar
orthonormal columns $v_i$.  Define
\[
 f_n=\sum_{i\le n}v_i,\qquad
 u_j=\sum_{i\le j}\frac{v_i}{j-i+1},\qquad
 A_1=0,\quad A_n=u_{n-1}u_{n-1}^T.
\]
The increments of $f_n$ are conditionally centered and spherical.
The weights are predictable and positive semidefinite; their norms
are bounded by $\pi^2/6$.  The manuscript checks the conditional
future-tail bound using the Hilbert-matrix norm bound $\pi$.
Nevertheless the embedding energy is
\[
 \sum_n\langle f_n,A_nf_n\rangle
 =\sum_{j=1}^{m-1}H_j^2,
 \qquad H_j=\sum_{i=1}^j\frac1i.
\]
The available variance budget is at most $4m$, while the normalized
gain is at least $\log^2(m/2)/16$.

This disproves a dimension-free consequence of the listed coarse
properties: conditional isotropy, positive predictable weights,
bounded expected tails, and the corresponding budget are not enough.
It is not an HSS counterexample.  The orthogonal family grows with
$m$, whereas the HSS path has a fixed query width.  The example also
shows why an abstract $p^2$ response-moment estimate cannot simply be
replaced by an $O(p)$ estimate.  Adding a decreasing quotient form does
not fix the abstract implication, since the constant form $I$ already
satisfies that condition.

\subsection{Frozen-forward gain can be arbitrarily large}

The conditional amplifier in \texttt{base:new:amplifier} is an actual
fixed-basis refit construction, but its conclusion is conditional.
At $k=1,q=4$, choose $M$ blocks and an odd
$n\ge4\pi^2M$, put
$c=2\cos(\pi/n)-1$ and $s=\sqrt{1-c^2}$, and alternate a run of
$n$ small column rotations with a release step.  The relevant
forward transfer is
\[
 \begin{pmatrix}c&0&s\\-s&0&c\\0&1&0\end{pmatrix}.
\]
Its eigenvalues are $-1,e^{i\pi/n},e^{-i\pi/n}$, so its $n$th power is
$-I$ while $c^n$ remains close to one.  After a fixed warm-up, each
release extracts a uniformly positive coefficient.  The
frozen-forward response satisfies $W(G)\ge M/16$ at height
$2+M(n+1)$, and continuity gives a positive-probability neighborhood
with response at least $M/32$.

The neighborhood probability is not uniform in $M$.  In fact the
fully averaged independent unit-source gain in the same construction
is bounded by $5/4$.  Choosing still slower rotations suppresses the
actual source excitation: the manuscript obtains a source-weighted
joint average at most $14.402\,\beta/M$ while the conditional unit gain
is at least $M/16$.  Thus a uniform bound after freezing the whole
forward stream is false; the unconditional physical theorem is not
disproved.  Near-protected covariance must be charged jointly with
the response rather than separated from it.

\subsection{Temporal parity and finite moment closure have limited scope}

Immediate fresh-frame parity kills a current odd cross term, not a
future weighted one.  For the temporal kernel
$K_{rs}=\mathbb E[T_rWT_s^T]$, the surviving delayed contribution for
$r<s$ involves future readouts after time $s$.  Exact rational
calculations in \texttt{base:new:temporal} and
\texttt{base:new:fourth}, using the original Haar law, exhibit nonzero
delayed entries.  For example, with the stated scalar $3$--$4$--$5$
rotations,
\[
 K_{35}=\frac{2925051264}{95367431640625}>0,\qquad
 K_{45}=\frac{4766750208}{476837158203125}>0.
\]
They are finite-family certificates, not an all-depth stability
theorem.  Fixed polynomial degree and a finite covariance state do
not give a spectral contraction; Haar integration also has signed
coefficients, so the monomial transfer is not a positive scalar
recursion.

Absolute-value path bounds lose another cancellation.  A sequence of
$h$ overlaps each of size $1/\sqrt{h+1}$ has bounded sum of squares
but an absolute sum of order $\sqrt h$.  Bessel control therefore
does not make the absolute ancestor sum bounded.  Orthogonality of
distinct initial sources removes cross-source terms but leaves
same-source temporal returns.

\section{Noncommutative refit estimates: the information that cannot be dropped}
\label{origobs:noncommutative}

\subsection{Source moments, covariance coercivity, and coherent rows}

The Gaussian/gamma moment-generating-function envelopes in
\texttt{base:new:source-moments} are inequalities for the original
source.  They do not identify its distribution or give stochastic or
convex order.  Likewise, auxiliary random signs used to take a square
root of a response are an analysis device, not replacement probes.
A tail bound for conditional expected risk is not automatically a
tail bound for realized reconstruction error.

For the exact source covariance and geometric loss,
\texttt{base:new:coercive} proves the sharp comparison
\[
 (1-\rho)^2\,\mathsf L\preceq\mathsf C\preceq\mathsf L,
 \qquad \rho=\frac{k}{q-2k}.
\]
At $q=4k$ the lower coefficient is $1/4$, attained by actual frames.
Thus an arbitrarily large improvement cannot come solely from
claiming that the exact covariance is much smaller than geometric
loss.  The comparison does not control how a dependent future
response aligns with that covariance.

The fourth-moment method in \texttt{base:con:gram} keeps
$\operatorname{tr}[(N^TN)^2]$ rather than replacing it by a crude
Frobenius scalar.  It gives a constant bound for suitably spread-out
source row energy, and isolates a bounded-dimensional exceptional
temporal eigenspace in general.  The hard case is a source with one
left row, $B=uv^T$, where $v$ still has all $mk$ temporal coordinates:
one source row does not reduce retained rank $k$ to one.  Padding that
preserves $B^TB$ does not change its spectrum; artificially spreading
the energy can cost a factor $k$.

An ideal Gaussian regression of the exceptional residual source
would pay the remaining trace.  It is not an available decoder:
the desired source is a component of $A-P_{\mathcal S}A$, while
sketches of $A$ also contain the unknown in-space part.  The
manuscript does not claim that this contaminating component can be
subtracted with the available information.  This is an access-model
obstruction to that proposed repair.

\subsection{Past output, the Haar fourth moment, and symmetric dilation}

The marked Bellman identity must contain both already emitted output
and the future value.  A terminal source mark reweights old
coefficients as well as future ones.  The exact telescope therefore
has local self terms and mixed terms involving the odd fresh-frame
kernel.  Improving the self-mark constant to
$8.257804885\ldots$ at width $4k$ does not pay those mixed terms
(\texttt{base:con:bellman}, \texttt{base:con:selfmark}).

Nor can all Haar fourth moments be majorized by their Gaussian
counterparts.  In the exterior-square sector the coefficients are
$2/[d(d-1)]$ and $2/d^2$, respectively; a determinant minor attains
the positive ratio $d/(d-1)$.  The valid marked analysis locates this
defect at the current mark rather than charging it through every
empty stage.  That placement does not eliminate the mixed
old/new-mark interaction.  A summable mean-forcing estimate likewise
cannot simply be added to a Bellman decomposition with which it
overlaps.

Symmetric dilation replaces $A$ by
$\left(\begin{smallmatrix}0&A\\A^T&0\end{smallmatrix}\right)$.
It doubles the rank and, with a common probe, changes the joint law
from the two independent-stream model.  Temporary diagonal-copy
coefficients must remain in the state.  In
\texttt{base:con:sym-cone}, the adapted state bound leaves an output
term of the form
\[
 \operatorname{tr}\bigl[L(PP^T-P^TP)\bigr],
\]
which has no fixed sign for nonscalar $L$.  A paid two-leg expression
$\|B^TKB\|_F$ does not dominate the required one-leg expression
$\|B^TKA\|_F$ without further hypotheses.  Scalar child weights make
the commutator vanish and yield an adapted bound, but do not authorize
substitution of an anticipative source.  Even the dilation of an
original rank-one problem has rank two, so it is not automatically
the scalar case.

\subsection{Refresh and observability must respect transported directions}

Matrix storage closes several restricted geometries: matched scalar
leakage, bounded exceptional levels, bounded runs between balanced
resets, and specified crossed noncommuting projector families.
These are sufficient geometric conditions, not properties of every
HSS learner.  Their observability condition must include transport.
For example, $A_0=e_1e_2^T$, $A_1=e_2e_1^T$ can have
$L_0+L_1=I$ while
$L_0+A_0L_1A_0^T=L_0$ and $\|A_0A_1\|=1$.  Positive untransported
window loss does not prevent indefinite survival.

Forcing strict refresh in every channel is incompatible with exact
recovery of a localized rank-one matrix $e_ie_j^T$ across opposite
root halves: its required direction must be preserved with norm one
at every ancestor.  Protected directions are not themselves an
instability.  If
$\Pi=\ker(I-F_hF_h^T)$, the actual source satisfies $N\Pi=0$
pathwise, and the manuscript bounds the corresponding root response.
The correct issue is source excitation of nearly protected
directions, not an unconditional requirement to contract every
direction (\texttt{base:new:protected}, \texttt{base:new:near}).

\section{Predictable feedback: false local charges and the later repair}
\label{origobs:predictable}

Several plausible predictable-energy shortcuts are disproved even
on reachable refit states.  In \texttt{centered:sec:obstruction},
a rank-one state with $g=e_1$, $w=e_2$, $Y=e_3$, $Z=0$, a row
identity completion and a column quarter-turn has unchanged quotient
energy, yet the current root Doob increment has variance
$1/(4d)>0$.  Therefore current predictable variance cannot always be
charged to current quotient dissipation.  Future reserves can still
pay it, so this is not a global risk counterexample.

The stronger positive-majorant obstruction in
\texttt{ch:rowmajorant} concerns replacing the exact quotient
$J_Y$ by a larger positive row energy $H_Y$.  Repeat the scalar
swap/identity completions
\[
 Q_U=\begin{pmatrix}0&-1\\1&0\end{pmatrix},\qquad Q_V=I.
\]
The deterministic feedback weights give unit weight to each adjacent
transpose innovation, while its actual transpose output weight
vanishes if another swap precedes the root.  On the reachable
canonical trajectory, $Y=0$ throughout and hence $J_Y=0$, but
$H_Y=C^2$ is positive.  Its scalar recursion tends to a nonzero
limit.  After the manuscript's three-step warm-up with one fixed
nonzero source block, the source budget is one and
\[
 \overline{\mathcal P}_{G,\mathrm{swap}}
   =\frac{n}{16d^2q}+O(1),\qquad
 \mathcal P_{G,\mathrm{swap}}=0.
\]
At $q=4$ the spurious slope is $1/256$.  Increasing to any other
fixed width does not remove the linear growth of this upper bound.
The lost terms in
$J_Y=H_Y-D_YD_Y^T-\Delta\Delta^T$ are needed for cancellation.

The earlier exact Doob decomposition, passive mean filter, five
future-variance registers, and bound on the marked state at each
deterministic time remain valid.  A current-state estimate is not an
estimate of a maximum or a sum over times.  Coordinate-averaged
response bounds leave a positive anisotropic remainder unless the
actual bases have the additional symmetry.  Backward weights are not
automatically monotone either.  These were missing premises in the
earlier proof route, not evidence that the final decoder fails.

\paragraph{Repaired status.}
Chapters \texttt{ch:feedback}--\texttt{ch:uniform} retain the future
projectors under legitimate frozen-stream conditioning, route
complementary loss, and preserve the signed retained-coefficient
combination.  In the notation of that proof, the cancellation
variable $R_i=f_ih_i^T+M_i^T$ has a transport recursion with the
otherwise duplicated input terms canceled.  Deterministic frame
transport then bounds the needed output, including coefficients
already emitted.  The resulting fixed-space theorem has factor
$1607$ with $8k$ fresh refit products; larger fixed widths yield
sharper constants and the stated near-optimal family.

Thus the historical logarithmic oversampling, exponential initial
refit bound, and unpaid predictable interaction are superseded as
obstacles to fixed-space refitting.  The counterexamples to their
discarded intermediate inequalities remain correct and explain why
the repaired proof keeps signed and noncommuting information.
None of these repairs establishes the quality of a noisy learned
space.

\section{Adaptive annihilation: generic erasure and actual inverse-growth families}
\label{origobs:actual-inverses}

The later chapters of \sourcefile{hss_complete_proof_chain.tex}
replace unrestricted adaptive examples by specified learners and fixed
input matrices.  The distinctions between their conclusions are
important: an erased sample, an exploding unweighted inverse, and a
positive physical approximation loss are three different outcomes.

\subsection{Perfect conditioning and comparator containment do not suffice}

The generic erasure construction in \texttt{ch:basisobs} fixes
$A=ue_j^T$, where $u=(u_a+u_b)/\sqrt2$ is supported in two children
and $j$ is external.  On a positive-probability event, the Gaussian
child maps have enough physical nullspace to choose unit vectors
$v_a,v_b$ satisfying
\[
 \Omega_a^Tv_a=a_0,\qquad \Omega_b^Tv_b=b_0,
\]
where $a_0$ is the normalized exterior row $e_j^T\Omega$ and $b_0$
is perpendicular to it.  The compressed design is a row isometry.
Nevertheless its row space contains the exterior anchor, so
nullification erases the outgoing rank-one observation.  The row
templates can contain the comparator exactly, and all these choices
can be nested.

This disproves an inference from containment and good conditioning
alone.  The selected column directions are admissible adaptive
choices, not claimed to be the actual canonical rule's trajectory.
The same chapter's path example shows a different limitation: a
fixed HSS-rank-one matrix can introduce a new cut direction at every
ancestor while the old direction restricts to zero in the next
child.  The rank-$k$ version has $kh$ such innovations along a path.
It excludes an argument that assumes only $k$ directions ever occur
along a path, but gives no $kh$ query lower bound because different
levels and nodes may be probed simultaneously.

\subsection{An actual phantom-direction family}

Theorem \texttt{actual:thm:inverse} concerns the canonical
un-nullified commutator learner.  Fix $B=e_jv^T$, where $v$ is a unit
vector supported in the opposite root half and is nonzero on every
coarse leaf.  On that half the true outgoing block is zero, but the
commutator contains
$-\zeta(v_I^T\Omega_I)$, with $\zeta=\Psi^Te_j$.
The full raw lift retains a direction related to
$(G^\dagger)^T\zeta$.  This unused, or phantom, direction attracts the
query images of both children toward the same exterior anchor.

For retained width $r>1$, leaf width $m=2r$, and $q>2r$, put
\[
 \mu=\frac{q-m}{m-2},\qquad
 \beta=\frac{q-2r+1}{q-r}<1.
\]
The exact graph-angle recursion and a Gaussian norm estimate imply,
at a parent whose children have height $h$ and
$m2^h\ge\max(q,8)$,
\[
 \Pr\!\left\{\sigma_{\min}(G^{(h+1)})^2
                   \le36m\mu\beta^h\right\}\ge\frac12,\qquad
 \mathbb E\|(G^{(h+1)})^\dagger\|_{\rm op}^2
                   \ge\frac{\beta^{-h}}{72m\mu}.
\]
The designs remain full row rank almost surely.  At $r=4,m=8,q=16$,
the lower bound is $(4/3)^h/768$.  Increasing to another fixed width
ratio slows the rate but does not restore a depth-uniform full
inverse moment.

The matrix is still represented exactly.  The collapsing orientation
has zero true outgoing comparator, and the nonzero row and column
ranges remain contained.  Thus this is an actual-rule obstruction to
an \emph{unweighted inverse intermediate bound}, not a lower bound
on approximation error.

\subsection{A saturated symmetric family removes the inactive-cut explanation}

Theorem \texttt{ports:thm:inverse} uses
\[
 B=\begin{pmatrix}0&UU^T\\UU^T&0\end{pmatrix},
\]
where $U$ has $k$ orthonormal columns and rank-$k$ restrictions on
every leaf; repeated, appropriately normalized identity blocks give
an explicit choice.  Every relevant incoming and outgoing cut has
rank $k$.  For $r\ge2k$ and $q\ge2r+k$, the raw commutator has
rank at most $2k$, so the actual full-range branch retains both the
signal and its external query port.  The children's query images
again align around that common port.

With
\[
 t_0=\frac{k(q-2r)}{2r-k-1},\qquad
 \lambda=\frac{q-2r+k}{q-r}<1,
\]
the exact angle recursion yields, at a specified height-$j$ node,
\[
 \mathbb E\|G_I^\dagger\|_{\rm op}^2
 \ge\frac{\lambda^{-(j-1)}}{8rq\,t_0}.
\]
At $r=4k,q=16k$ this is
$(7k-1)(4/3)^{j-1}/(4096k^4)$.  The proof uses
Cauchy--Schwarz and Jensen on the joint angle and physical-probe
norm; it does not factor them as independent.  Expectations can be
extended nonnegative expectations; finiteness at each height is not
part of the conclusion.

Exact representation still holds.  In particular the error
projector annihilates the lifted external port:
\[
 (I-U_I^{\rm loc}(U_I^{\rm loc})^T)(G_I^\dagger)^TH=0.
\]
Even saturation and nonzero cut signals therefore do not justify a
uniform full inverse bound.  The correct error readout may cancel
the very directions responsible for inverse growth.

\section{Inherited error and the actual noisy annihilator failure}
\label{origobs:annihilator-failure}

\subsection{A clean parent can create new loss with well-conditioned designs}

Theorem \texttt{return:thm:example} follows the actual hybrid rule
at $k=1,r=2,m=4,q=16,N=32$.  For adjacent leaves $a,b$ and an
external coordinate $j$, set
\[
 B=(e_{a,1}+e_{b,1})e_j^T,\qquad
 E[a,b]=e_3e_1^T+2e_2e_2^T+\tfrac32e_4e_3^T.
\]
All other blocks of $E$ vanish.  Thus $E$ is internal to the parent
$I=a\cup b$, and its current cut commutator is identically zero.
On explicitly specified open probe events the current row design
has singular values in $[3/4,5/4]$, the column design is close to a
row isometry, and the surviving external anchor is bounded away
from zero.  The current raw sample tails vanish in both orientations.

Nevertheless, for $A_\varepsilon=B+\varepsilon E$, the inherited
and newly discarded row losses at the reference probes are
\[
 h_U^2=\frac{\varepsilon^2}{1+\varepsilon^2},\qquad
 e_U^2=
 \frac{\varepsilon^2M^2(1+\varepsilon)^2}
 {((1+\varepsilon)^2+\varepsilon^2)^2+
                  \varepsilon^2M^2(1+\varepsilon)^2}.
\]
Here $M$ is arbitrary and
$0<\varepsilon\le\min(0.1,0.1/M)$.  Consequently
$e_U^2/h_U^2\to M^2$ as $\varepsilon\downarrow0$.
Strict local singular-value gaps make the inequalities persist on
positive-probability neighborhoods.

This excludes charging new physical loss only to the current
crossing residual or current truncation tails.  It also excludes a
uniform pointwise inherited-loss gain under those current-design
safeguards.  The event probabilities need not be uniform in $M$,
so the result is compatible with the manuscript's unconditional
mean estimates in the narrower clean model.

The failure occurs despite a valid all-depth theorem for the
unlifted data.  If $E=A-B$ for any fixed HSS-rank-$k$ comparator,
Theorem \texttt{actual:thm:datafit} proves for the actual hybrid
selections, including singular internal designs,
\[
 \sum_I\mathbb E\bigl(
       \|\mathcal R_{U,I}\|_F^2+\|\mathcal R_{V,I}\|_F^2\bigr)
       \le6q^2L\|E\|_F^2.
\]
The proof uses only that each projected comparator commutator has
rank at most $2k$ and that the selected query space minimizes
the corresponding rank-$2k$ tail.  Its first physical leaf
layer also has a residual-relative bound.  Neither result
converts a small query-space residual into a small physical
loss under an adaptive opposite annihilator.

There is an additional useful distinction between a vanishing
\emph{local} tail and vanishing tails throughout the tree.
Theorem \texttt{return:thm:zero} proves, for each fixed $A$ and
the stated hierarchical branch, that the total tail is zero
almost surely exactly when $A$ belongs to the HSS-rank-$2k$
class; the learned space then contains $A$.  The probability-one
event may depend on $A$ and is not simultaneous over all
sketch-dependent matrices.  This qualitative endpoint is
compatible with arbitrarily poor quantitative behavior near
it and with the local clean-parent example.

\subsection{A fixed-input family disproves the annihilate-first basis bound}

Theorem \texttt{anchor:thm:main} is stronger: it disproves the desired
basis-quality bound for the specified unchanged annihilate-first
hybrid.  Fix $r\ge2$, $q>4r$, and a depth parameter $h$; put
$n=4r2^h$, $N=4n$.  Let $I$ and $K$ be adjacent blocks of size
$n$, $P=I\cup K$, and $j$ the first coordinate outside $P$.
Set $u=n^{-1/2}\mathbf1_I$ and $B_h=ue_j^T$.  Choose two unit
Frobenius-norm perturbations with mutually disjoint supports:
\[
 \begin{aligned}
 E_0&=\frac1{\sqrt n}
       \{\text{permutation exchanging adjacent $2r$-leaf blocks in $I$}\},\\
 E_1&=\frac1{\sqrt r}\sum_{a=1}^r e_{n+a-1}e_{j+a}^T.
 \end{aligned}
\]
The fixed input is
$A_h=B_h+\delta_hE_0+\varepsilon_hE_1$.
The comparator is HSS rank one and
$e_k(A_h)\le\delta_h^2+\varepsilon_h^2$.

The arbitrarily small positive perturbation $\delta_hE_0$ activates
otherwise unused column directions.  Across the clean subtree these
directions trap the row anchor
$g=e_j^T\Omega$ in the learned column query spaces.  The original
probe reuse is essential to this accumulation.  With
\[
 \alpha=\frac{q-2r+1}{2(q-r)}<\frac12,\qquad
 \mu=\frac{q-2r}{2r-2},
\]
the expected squared surviving anchor is at most
$q\mu\alpha^h$.  The sibling perturbation
$\varepsilon_hE_1$, with $\varepsilon_h=\alpha^{h/4}$, then
displaces the almost annihilated signal in the rank-$r$ truncation.
At $\delta=0$ the canonical zero-sample branch is different; exact
recovery on that stratum gives no continuity through positive
$\delta$.

For each fixed $h$, bounded convergence permits a
\emph{deterministic} choice
$0<\delta_h\le\varepsilon_h^2$ along a predetermined sequence.
It is chosen by an expectation criterion, never from the realized
probes.  In the maintained parameters $r=2k,q=16k$, the theorem is
\[
 \mathbb E\,d_{\mathcal S_{\rm hyb}(A_h)}(A_h)^2\longrightarrow1,
 \qquad
 L_he_k(A_h)\longrightarrow0,
 \qquad L_h=h+4.
\]
The optimum is strictly positive for every $\delta_h>0$, since an
interior sibling block has rank $2r>k$.  Therefore no finite
constant $C_b$ can make
$\mathbb E d_{\mathcal S_{\rm hyb}(A)}(A)^2\le C_bLe_k(A)$
hold for all fixed inputs.

The first adaptive-merge theorem is not contradicted.  There the
child templates have a stream-separation property.  Here a
zero-comparator orientation can store harmful query information at
zero immediate comparator cost; at later levels it acts through
the opposite annihilator, where stream separation is absent.

\subsection{What the counterexample says about fallbacks and unions}

Normalized direct sums of $T$ aligned copies preserve HSS rank one
while making the comparator have global rank $T$ with equal
singular values $T^{-1/2}$.  The expected hybrid-space loss still
tends to one.  For $T>4k$, any global-rank-$4k$ candidate has limiting
squared error at least $1-4k/T$, and a physical-leaf-diagonal
candidate misses the entire unit comparator.  Hence selection among
these fallbacks before further nonlinear HSS projection does not
repair the example
(\texttt{anchor:cor:directsum}, \texttt{anchor:eq:globalfallback}).
No assertion is made about arbitrary subsequent best-HSS
projection, which can change the bases.

The un-nullified companion retains the displaced signal on this
family.  The original manuscript proves exact two-view rescue and
then a mean-square small-noise result in
\texttt{ch:tangent} and \texttt{ch:nonlinear}.  These are positive
results for the specified family and interfaces.  The two complete
hierarchies are learned separately before their physical union is
taken.  They are not a recursively union-fed learner, and they can
have different child templates.  A proof about a distinguished
clean cut is not a whole-space, all-input basis theorem.

In particular, the two-view hyperplane theorem conditions on a
full-column-rank lift, a projector, and a nonzero anchor in both
projector sectors, while leaving an independent Gaussian
competitor block.  Its projector-complement rank must be at
least $r+1$.  The two leading-$r$ spaces then span the full
$(r+1)$-dimensional lift almost surely
(\texttt{tangent:thm:twoview}).  Learned opposite designs on an
arbitrary input cannot simply be declared independent to fit
this theorem.

The tangent result first fixes height and lets the perturbation
vanish.  The nonlinear moment argument supplies the stated
expectation interchange in that local model.  It does not prove a
height-uniform finite-noise remainder.  Changing the order of the
height and perturbation limits, or importing the same subtree
locality to general coupled HSS cuts, would be an additional premise.

\section{Rank slack and simultaneous interpolation do not resolve the interface}
\label{origobs:rank-slack}

For $S=C+F$ with $\operatorname{rank}C\le k$, the leading-$r$
projector $P$ satisfies the valid local spectral estimate
\[
 \|(I-P)S\|_{\rm op}^2
       \le\frac{\|F\|_F^2}{r-k+1}.
\]
Consequently, if $\operatorname{rank}Z\le k$ and
$\|Z\|_{\rm op}\le1$, then
$\|(I-P)SZ\|_F^2\le k\|F\|_F^2/(r-k+1)$.
The target here is the observed sum $S$, not its comparator part
$C$.  Taking $S=0$, $F=-C$, with $C$ perpendicular to the
canonical zero-sample completion, prevents replacing $S$ by $C$
with a universal contraction coefficient below one
(\texttt{ch:slack}).

Likewise, capturing both local innovations does not necessarily
produce a matrix interpolating both complete sketches.  In
\texttt{ch:bicriteria}, the two-sided construction can annihilate
the forward residual yet leave a transpose term
\[
 (I-Q)(B_Z^T-B_Y^T)(I-P)\Psi.
\]
The underlying generic dimension obstruction is explicit.  At
width $q=3r$, the two-sketch measurement rank is
$6Nr-9r^2$, whereas the saturated HSS-rank-$r$ parameter count
in the stated tree model is $5Nr-6r^2$; the former is larger for
$N>3r$.  Generic observations therefore have no exact
HSS-rank-$r$ interpolant.  This rules out an exact interpolation
step in that argument, not an approximate HSS recovery theorem.

Finally, adding a large physical-leaf-diagonal matrix does not
regularize a raw learning trajectory: on a full-rank trajectory,
$D+\varepsilon E$ has the same learned bases as $E$, and the
homogeneous reconstruction is $D+\varepsilon\widehat E$.
Thus closeness to an exactly representable leaf-diagonal input is
not, by itself, a stability proof (\texttt{ch:fallback}).

\section{Local equivariance is insufficient; guards leave a mandatory cost}
\label{origobs:guards}

\subsection{A broad local-equivariant model has genuine physical growth}

Chapter \texttt{ch:budget} fixes leaf sources before one Gaussian
probe and permits subtree-local, query-rotation-equivariant
selections.  It proves the exact innovation and energy identities
before constructing a bad selection.  Write
\[
 \lambda=\frac{q-2r}{q-r},\qquad
 \kappa=\frac{q-r-1}{q-2r-1},\qquad
 \rho_0=\frac r{q-r-1}.
\]
For child height $t$, the normalized innovation energy is
$m_I=\lambda^t\beta_I$.  If $V_I$ is the retained inverse trace
and $a_I$ its innovation-weighted mean, the merge cost is
\[
 \mathbb E\|z_I\|_F^2
 =(\kappa-1)(a_a+a_b)+\kappa(m_av_b+m_bv_a),
 \qquad v_i=\mathbb EV_i.
\]
The cross-child terms already show why innovation decay alone
does not control physical error.

In Theorem \texttt{budget:thm:counter}, one root subtree has a
single unit source and fixed coordinate selections; the other
has zero source but retains the weakest $r$ Gram directions at
every merge.  Its errors vanish, while its inverse mean grows.
At the root an admissible selection discards the entire fitted
coefficient vector.  All designs are full row rank and all
second moments are finite at each finite height, but
\[
 \mathbb E\|E_{\rm root}^{\rm out}\|_F^2
 \ge1+\rho_0+\kappa\rho_0(\lambda\kappa)^h,\qquad
 \lambda\kappa
 =1+\frac{r}{(q-r)(q-2r-1)}>1.
\]
At $r=2,q=16$ the lower bound is
$15/13+(2/11)(78/77)^h$.
For $r=1$ the sharper sphere-angle calculation replaces $\kappa$
by
\[
 \chi_q=\frac{q-2+2C_q}{q-3},\qquad
 C_q=\frac{\Gamma(q/2)}
            {\sqrt\pi\,\Gamma((q-1)/2)},
\]
and gives $\lambda\chi_8>1.3778$.

This disproves a theorem uniform over the entire local-equivariant
selection class.  The construction deliberately chooses weak
directions and, at the root, discards the fitted target.  It is
not the canonical full-fit rule or a fixed-matrix counterexample
to the maintained HSS union.  Its purpose is to show that locality
and rotational symmetry do not replace analysis of the actual
selection.

\subsection{A sufficient relative guard can conflict with exact containment}

The same chapter proves a positive conditional theorem under
$V_I\le\gamma\min(V_a,V_b)$ at ordinary nodes.  The physical
bound contains $\sum_{t<H}(\gamma\lambda)^t$; it is uniform in
height when $\gamma\lambda<1$ and linear at equality.  Choosing
the leading Gram directions supplies $\gamma=1$, but can discard
a required signal.  Stability of this selection therefore does
not by itself define a valid HSS learner.

Proposition \texttt{budget:prop:infeasible} makes the conflict
quantitative.  Let
\[
 K=\begin{pmatrix}
 I_r&(1-\varepsilon)e_1e_1^T\\
 (1-\varepsilon)e_1e_1^T&I_r
 \end{pmatrix},\qquad
 Q=\frac1{\sqrt2}\binom{e_1}{-e_1}.
\]
Both child inverse traces equal $r$, but the mandatory direction
$Q$ has Gram eigenvalue $\varepsilon$.  Among all rank-$r$
compressions containing it,
\[
 V_{\rm opt}=\frac1\varepsilon+\frac1{2-\varepsilon}+r-2.
\]
No fixed finite relative guard can preserve every prescribed
signal on every positive-definite realization.  This is a local
algebraic obstruction; it says nothing about the frequency or
source weight of the bad configuration under a specific learner.

\subsection{Optimal padding isolates, but does not remove, the core obligation}

The exact optimizer in \texttt{budget:thm:completion} retains an
orthonormal mandatory frame $Q$ of rank $s\le r$.  In coordinates
$[Q,Z]$, write
\[
 K=\begin{pmatrix}A&B\\B^T&D\end{pmatrix},\qquad
 \mathcal S=D-B^TA^{-1}B,\qquad
 \mathcal M=I+B^TA^{-2}B,\qquad
 \phi=\operatorname{tr}(A^{-1}).
\]
If $\nu_j$ are the increasing eigenvalues of
$\mathcal S^{-1/2}\mathcal M\mathcal S^{-1/2}$, block inversion
and generalized Rayleigh--Ritz give
$V_{\rm opt}=\phi+\sum_{j=1}^{r-s}\nu_j$.
This can be computed locally without new matrix products.
The later additive theorem \texttt{core:thm:additive} sharpens
the bookkeeping to
\[
 V_{\rm opt}\le
 \phi+\sum_{i=s+1}^{r}\lambda_i(K)^{-1}
 \le\phi+\min(V_a,V_b),
\]
with the eigenvalues of $K$ in decreasing order.  The unavoidable
quantity is the mandatory-core inverse $\phi$, not optional
padding.  A full-rank truncated sample has $s=r$ and no optional
freedom at all.

The fourth-moment result \texttt{core:thm:fourth} is valid under
the local model even for unstable admissible selections:
\[
 \mathbb E\|w_I\|_F^4
 \le d_0(d_0+2)\chi^t\beta_I^2,\qquad
 d_0=q-r,\quad d=q-2r,\quad
 \chi=\frac{d(d+2)}{d_0(d_0+2)}.
\]
Thus a uniform marginal $L^2$ bound on $\phi$ pays its dependent
product with the innovation by Cauchy--Schwarz.  At $r=2,q=16$,
the decay is $\sqrt\chi=\sqrt3/2$, not $\lambda=6/7$.
The additive inverse-budget theorem allows inverse means to grow
linearly while the physical cost remains summable.

For a \emph{fixed physical} nested core, the core projection has
the required Gaussian marginal law even if its optional
completion is adaptive.  The original chapter obtains explicit
uniform constants for this restricted setting.  A noisy selected
core need not be fixed; replacing it by an unknown comparator
changes the learner and invalidates that law.  General cuts can
also share exterior probe rows and use unequal opposite
templates, violating the independent-subtree interface.  These
are the remaining premises identified in
\texttt{core:sec:pilot} and \texttt{core:sec:checks}, not
counterexamples to every possible guarded algorithm.

\section{Coverage of the historical ledger and scope of the remaining premise}
\label{origobs:coverage}

The following map accounts for all 66 numbered routes in the
original \nolinkurl{base:new:route-history}.  Numbers are grouped by
mathematical issue rather than by the order in which they were
investigated.  Positive restricted theorems appear because their
hypotheses delimit the corresponding obstruction.

\begin{center}
\small
\begin{tabular}{p{.18\linewidth}p{.75\linewidth}}
\hline
Original routes & Taxonomy and precise qualification\\
\hline
1, 2, 7 &
Section \ref{origobs:interfaces}: exact recovery and the first two
levels do not imply all-depth reuse; freshness changes query cost.\\
5, 6, 10 &
Section \ref{origobs:adaptive-inverses}: adaptive inverse moments,
Stein correction, and legitimate conditional Wishart laws.\\
8, 9, 12, 55, 56 &
Section \ref{origobs:validation}: tails, validation, fitting
injectivity, saturated identification, and local versus uniform
conditioning.\\
3, 4, 14, 15, 17, 30 &
Section \ref{origobs:quotient}: inactive reconstruction, common
coefficients, signed merges, correlated feedback, and scalar storage.\\
13, 16, 18 &
Sections \ref{origobs:embedding} and \ref{origobs:coverage}:
finite-profile hypotheses, restricted coordinate paths, and the
loss in absolute overlap bounds.\\
19--29, 36--38, 45, 46 &
Section \ref{origobs:embedding}: source orthogonality, full covariance
state, conditional amplification, temporal marks, embedding,
near-protection, moment envelopes, and repaired restart/source
accounting.\\
31--35, 39--44, 47--54, 57, 58 &
Section \ref{origobs:noncommutative}: directional storage,
transported refresh, covariance metrics, coherent rows, inaccessible
auxiliary data, Bellman marks, Haar defect, dilation, and symmetry.\\
11, 59--66 &
Section \ref{origobs:predictable}: initial refit bounds, passive
filter, Doob product, quotient/variance interfaces, current-state
control, and their later fixed-space repair.\\
\hline
\end{tabular}
\end{center}

For completeness, the finite-profile transfer theorem called
$F^+$ in \texttt{base:new:fp} assumes a specified centered
rank-one conditional response profile.  Isotropy of an HSS sketch
does not imply that profile.  The factor-$10$ coordinate-path
theorem exploits a nonbranching coordinate geometry.  Neither is
a substitute for proving the hypotheses along a general adaptive
trajectory.  The source-history ledger's finite temporal
certificates and repeated-eigenvalue symmetry arguments have the
same restriction: their finite family or actual symmetry must
be present.

Several ideas were considered but never promoted to completed
algorithms: repeated independent residual corrections, redundant
retained dimensions intended to balance arbitrary bases,
query-coordinate gauge randomization, and a direct whole-space
small-ball estimate.  The record contains no theorem showing
that one of them simultaneously preserves exactness on the
original rank class, costs $O(k)$ products, and has a uniform
residual-relative bound.  Their absence is a missing construction,
not an impossibility theorem.

The later obstruction families, not contained in the first
66-route chronology, are covered explicitly above: independent
half solves, adaptive erasure and path novelty, tree-nuclear
selection, rank-slack and interpolation limitations, both actual
inverse-growth families, the well-conditioned inherited-return
example, the actual annihilator failure and its fallbacks,
the broad local-equivariant counterexample, and mandatory-core
guard infeasibility.  All statements in this chapter concern the
rules, models, and probability quantifiers stated with them.

The original manuscript therefore ends with a sharply separated
dependency chain.  Independent refitting of a fixed nested space
is proved.  Exact offline rank reduction is a valid query-only
step.  The specified annihilate-first learner's general basis
bound is false.  The maintained two-hierarchy union has exact and
small-noise results on the stated special family, but its
all-input basis-quality estimate is not supplied by those
results or by the abstract local regression theorems.  This is
the original manuscript's remaining interface, not a lower
bound against all $O(k)$-query HSS algorithms.

\paragraph{Verification scope.}
The exact fractions quoted above are historical finite
certificates preserved in the original manuscript.  No
Monte Carlo experiment or historical certificate engine was
rerun for this consolidation.  A finite simulation can check
an identity or illustrate a specified trajectory; it does not
establish an all-depth probabilistic estimate.  All rank,
inverse, and nullspace statements here concern exact
arithmetic with the displayed conditioning, rather than a
floating-point rank threshold.

\chapter{Obstructions and limits established after the original manuscript}
\label{chap:later-obstructions}

This chapter records what the subsequent investigations ruled out, what
they repaired, and what remained unresolved when research stopped.  Its
scope is deliberately narrower than an impossibility theorem for HSS
approximation.  A counterexample to a selector, an adaptive moment
estimate, or one prescribed hierarchy is not a lower bound against every
algorithm using a depth-independent number of queries.  The general
conjecture---an $O(k)$-query algorithm with expected squared Frobenius
error $O(L\operatorname{OPT}_k(A))$ and final HSS rank at most $k$---was
not proved or disproved.

All matrices in the probabilistic counterexamples below are fixed before
the Gaussian probes.  Some competing interpolants, conditional states,
or event windows are subsequently selected from the observations; those
quantifiers are stated explicitly.  The original manuscript's adaptive
nullification and annihilate-first counterexamples are recalled only
where needed to distinguish later learners.  The detailed follow-up
catalogue is based on the source notes and their recorded audits, not
on the numerical summaries alone.  \sourcefile{hss_audit_report.md}
provides the chronological overview.

\section{Global interpolation and minimum-norm selectors}
\label{sec:later-global-selectors}

\subsection{A fixed rank-one input with smaller rank-one HSS interpolants}

The global-selector extension removes the qualification that a
regularizer obstruction requires nominal rank at least two.  Let the
two root children $I,J$ each have size $n$, let $q<n$, and draw
independent standard Gaussian probes $\Omega,\Psi\in\mathbb R^{2n\times q}$.
For fixed unit vectors $u,v\in\mathbb R^n$, set
\[
 B=\begin{bmatrix}0&uv^T\\0&0\end{bmatrix}.
\]
This fixed input is HSS rank one on the complete coordinate tree and has
unit Frobenius norm.  Almost surely the systems
\[
 \Psi_J^Ta=\Psi_I^Tu,\qquad
 \Omega_I^Tb=\Omega_J^Tv
\]
have solutions; choose their minimum-Euclidean-norm solutions.  For every
real $t$, the observed-data-dependent competitor
\[
 C(t)=\begin{bmatrix}
 (1-t)ub^T&tuv^T\\
 -(1-t)ab^T&(1-t)av^T
 \end{bmatrix}
\]
satisfies
\[
 C(t)\Omega=B\Omega,\qquad C(t)^T\Psi=B^T\Psi.
\]
It remains HSS rank one: at any node contained in $I$, outgoing and
incoming exterior blocks have the respective common factors $u_S$ and
$b_S^T$; at a node contained in $J$, the corresponding factors are $a_S$
and $v_S^T$.  Thus the assertion concerns every cut, not merely the root
split.

Writing
\[
 \eta=\|a\|^2+\|b\|^2+\|a\|^2\|b\|^2,
\]
the disjoint supports of the four blocks give
\[
 \|C(t)\|_F^2=t^2+(1-t)^2\eta.
\]
Its minimum occurs at $t_*={\eta}/({1+\eta})$, and
\[
 \|C(t_*)\|_F^2=\frac{\eta}{1+\eta}<1,
 \qquad
 \|B-C(t_*)\|_F^2=\frac1{1+\eta}.
\]
Consequently minimum-Frobenius-norm HSS interpolation, including minimum
norm as a tie rule for exact empirical least squares, fails almost
surely on this fixed exact input.  A universal relative guarantee would
require zero error because $\operatorname{OPT}_k(B)=0$ for every $k\ge1$.

The exact displayed error belongs to the constructed competitor, not
to an arbitrary optimizer.  The latter nevertheless has a quantitative
risk lower bound.  If $n>q+1$, put
\[
 \rho=\frac{q}{n-q-1},\qquad \mu=2\rho+\rho^2.
\]
Inverse-Gram averaging gives $\mathbb E\|a\|^2=
\mathbb E\|b\|^2=\rho$; the two vectors use independent probe streams,
so $\mathbb E\eta=\mu$.  Concavity of $x/(1+x)$ and the reverse triangle
inequality in the space of square-integrable matrices imply, for every
measurable minimum-norm selection $\widehat C$,
\[
 \mathbb E\|B-\widehat C\|_F^2
 \ge\left(1-\sqrt{\frac{\mu}{1+\mu}}\right)^2.
\]
This tends to one as $n/q\to\infty$.  Existence of the minimizer follows
from closed rank constraints and compact norm sublevel sets.  These
statements are established in \sourcefile{hss_global_selector_obstruction.tex}.

\subsection{Tree penalties and exact leaf-block calibration}

For constant $u,v$ on balanced root halves, the outgoing-plus-incoming
cut nuclear norm of $B$ at a depth with $m$ nodes per half is $2\sqrt m$.
All cuts of $C(t)$ have rank at most one.  The partition bound
\[
 \sum_S\bigl(\|C[S,S^c]\|_F^2+
                    \|C[S^c,S]\|_F^2\bigr)\le2\|C\|_F^2
\]
and Cauchy--Schwarz bound its nuclear sum by
$2\sqrt{2m}\|C\|_F$.  On $\eta<1$, the smaller interpolant has
$\|C(t_*)\|_F<1/\sqrt2$ and improves every depth strictly.  It therefore
improves every nonzero nonnegative deterministic depth-weighted sum of
these cut nuclear norms.  Its ordinary matrix rank is at most two, so
its ordinary nuclear norm is also smaller than that of $B$ on this
event.  Markov's inequality gives probability at least
$\max(0,1-\mu)$ for the event.  The same strict comparison defeats
positive penalization because both empirical residuals vanish.

Supplying all physical-leaf diagonal blocks exactly does not repair
these selectors.  Mark one leaf in each sibling pair active and the
other inactive.  Support $u,v$ on the active leaves of their respective
halves, and restrict $b,a$ to inactive leaves.  The free dimension in
each interpolation system is then $d=n/2$; assume $d\ge q$.  On every
physical leaf $K$, at least one factor of the corresponding diagonal
outer product vanishes, so
\[
 C(t)[K,K]=B[K,K]=0
\]
exactly.  The smaller-norm construction is otherwise unchanged.

For $d>q+1$, replace $\rho,\mu$ by
$\rho_d=q/(d-q-1)$ and $\mu_d=2\rho_d+\rho_d^2$.  The minimum-norm
lower bound remains valid.  With $u,v$ spread uniformly over active
coordinates, the true leaf-depth nuclear sum is $2\sqrt{m/2}$, while
the competitor still obeys $2\sqrt{2m}\|C\|_F$.  The stronger event
$\eta<1/3$, of probability at least $\max(0,1-3\mu_d)$, gives
$\|C(t_*)\|_F<1/2$ and improves every depth, including the leaves.
Any finite penalty depending only on the supplied leaf diagonal blocks
takes the same value on both candidates and cannot reverse this
comparison.  A single-active-leaf version also gives calibrated
minimum-norm failure, with free dimension $n-s$ for leaf size $s$.
See \sourcefile{selector_with_leaf_data.md} and the final theorem in
\sourcefile{hss_global_selector_obstruction.tex}.

The alternatives above depend on the probes.  They disprove specified
global selection principles; they do not say that two fixed inputs are
indistinguishable with positive probability, or that a different
structured decoder cannot recover each fixed HSS input exactly.

\subsection{Certificates restore exact uniqueness, not robust selection}

The trajectory-certified selector imposes injectivity of each complete
physical template together with its exterior cut anchor.  Under these
certificates, equal transcripts determine equal outgoing and incoming
intersections, hence equal canonically completed child frames.  Induction
then gives a common coefficient space; successive two-sided compression,
ending with an injective root core, proves exact uniqueness.  The false
interpolants above violate these full-template certificates, including
the endpoint where a minimal cut anchor alone would not detect them.

For a fixed exact input the comparison templates are deterministic, so
each $q\times p_I$ certificate matrix is Gaussian and
$\mathbb E\sigma_{\min}^{-2}\le p_I/(q-p_I-1)$.  This marginal
bound cannot be replaced by a tree-size-independent bound on the
maximum certificate cost: disjoint leaves provide arbitrarily many
independent Gaussian direction images.  Hard thresholds also exclude
some exact inputs with positive probability.  The quantitative
compression inequality gives products of levelwise factors, rather
than the desired robust residual-relative bound for changing candidate
spaces.  The exact measurable selector is therefore a valid positive
endpoint, not a completed approximation algorithm.
See \sourcefile{certified_selector_route.md} and
\sourcefile{certified_selector_audit.md}.

\section{The one-sided Gaussian pilot obstruction}

Removing transpose data from a pilot would make its row spaces
independent of a later transpose stream, but exact bounded-rank row
containment already fails for that restricted acquisition model.
Fix query width $q$ and output rank $r$.  Take powers of two
$d\ge2r$ and $M\ge2q+2$.  Partition one root half $I$ into $d$ aligned
subtrees $K_i$ of size $M$, choose fixed unit coordinates $e_i\in K_i$,
and choose a fixed unit coordinate $v$ in the other half.  For arbitrary
$t_i$ and $b_i$ supported in $K_i$, define
\[
 u=\sum_{i=1}^d t_i e_i,
 \qquad A=\sum_{i=1}^d e_i b_i^T+uv^T.
\]
Every such matrix is HSS rank one, including at cuts inside the $K_i$.
The all-cut verification uses the single row factor $e_i$ inside a
block and the single exterior factor $v$ above unions of complete
blocks; no large unrestricted leaf is hidden in the construction.

A pilot observing only $(\Omega,A\Omega)$ obtains, apart from zero
rows, exactly
\[
 y_i=X_i^Tb_i+t_i\gamma,
 \quad X_i=\Omega_{K_i},\quad \gamma=\Omega_J^Tv.
\]
To obtain a fixed-input impossibility, use an input prior only in the
proof: independently take $t_i\sim N(0,1)$ and
$b_i\sim N(0,\tau^2 I_M)$ before drawing $\Omega$.  Conditional on
the complete observations, the posterior covariance of $t$ is diagonal
with strictly positive entries
\[
 w_i=\left[1+\gamma^T(\tau^2X_i^TX_i)^{-1}\gamma\right]^{-1}.
\]
Thus $u$ has a posterior density on the entire fixed $d$-dimensional
span of the $e_i$.  After also fixing any independent algorithmic seed,
a returned space of dimension at most $r<d$ contains $u$ with posterior
probability zero.  Fubini converts this statement into failure for
almost every fixed input under the proof prior, with the bad-input set
allowed to depend on the pilot.

There is also a quantitative Bayes loss bound.  A rank-$r$ projector
captures at most $r$ from a covariance bounded by the identity, giving
\[
 \mathbb E\|(I-P_S)u\|^2
 \ge\frac{d}{1+q/[\tau^2(M-q-1)]}-r.
\]
For $\tau=2$ and the displayed dimensions this is at least $3d/10$.
The result permits unlimited computation and disclosure of the entire
parametric family.  It does not cover structured forward queries: on
this very family the single query $Av=u$ reveals the desired vector.
Transpose products and independently trained two-sided pilots are also
outside its scope.  The theorem and independent review are
\sourcefile{one_sided_gaussian_pilot_obstruction.md} and
\sourcefile{one_sided_gaussian_pilot_audit.md}.

\section{Why the screened-source moment proposed by Claude is false}

\subsection{The valid deterministic cancellation}

For a full-row-rank design $G=U\Sigma V^T$, lift
$L=(G^\dagger)^T$, response $J$, and retained projector $P$ with
$PLJ=LJ$, split the singular directions into the weakest $d$ and the
remaining directions.  If $J_1=V_1^TJ$ has full row rank, then
\[
 (I-P)L=(I-P)U_2\Sigma_2^{-1}
       \bigl[V_2^T-J_2J_1^\dagger V_1^T\bigr].
\]
This algebraic identity is correct.  The failed proposal was to discard
the remaining left projector and assert a uniformly controlled second
moment for the resulting screened source.  Weak singular values must
indeed be the removed block; sorting in decreasing order would not even
give the claimed deterministic inverse bound.  Full rank of $G$ also
does not imply the required response rank of $J_1$.

\subsection{An infinite moment at an ordinary Gaussian leaf}

The counterexample uses the actual full-range learner's leaf dimensions
$m=16,q=64,s=4,r=8$.  Fix orthonormal physical vectors $a,e$ inside a
leaf and a unit $v$ outside it, and take
\[
 A=av^T,\quad B=(a-e)v^T,\quad E=ev^T.
\]
The comparator is legal but need not be best.  Here $G$ is an ordinary
Gaussian leaf design and, for an independent Gaussian response row
$\gamma=v^T\Omega$,
$J=G^Ta\gamma$ and $J(E)=G^Te\gamma$.
Sort the singular values increasingly, and write
$c=U^Ta,d=U^Te$.  Screening only the weakest coordinate gives
\[
 Q(E)_i=\sigma_i(d_i-c_i d_1/c_1)\gamma,
 \qquad i\ge2.
\]
Consequently
\[
 \sigma_2^{-2}\|Q(E)\|_F^2
 \ge\|\gamma\|^2\left(1+\frac{d_1^2}{c_1^2}\right).
\]
Haar orientation gives a density for $c_1$ positive at zero and
$\mathbb E[d_1^2\mid c_1]=(1-c_1^2)/(m-1)$.  The expectation is
infinite, even conditional on any nonzero $\gamma$.  All stated rank
assumptions hold almost surely.  In contrast, the actual fit is
$a\gamma$, so the retained space contains $a$ and
\[
 \|(I-P)LJ(E)\|_F^2\le\|\gamma\|^2.
\]
Its expectation is at most four.  The discarded left projection can
cancel a nonintegrable response ratio.

\subsection{Linear scaling even for an exact best comparator}

Fix orthonormal $a,e$ inside the leaf and orthonormal $v,w$ outside it.
For $0<\varepsilon<1$, let
\[
 A=av^T+\varepsilon ew^T,\quad B=av^T.
\]
The outgoing cut has singular values $1,\varepsilon$, and $B$ attains
the cut lower bound globally.  Thus
$\operatorname{OPT}_1(A)=\varepsilon^2$.  Put
\[
 \gamma=v^T\Omega,\quad \eta=w^T\Omega,\quad
 g=\|\gamma\|^2,\quad \kappa=\|\eta\|^2,\quad
 h=\gamma\eta^T,\quad \Delta=g\kappa-h^2.
\]
The same one-coordinate screening has a lower-bound integrand whose
Haar average is exactly
\[
 \frac{\varepsilon g\kappa}{2\sqrt\Delta}.
\]
One obtains this by setting $t=c_1/d_1$, using its standard Cauchy
density, and integrating the resulting rational function.  For
independent Gaussian response rows in $\mathbb R^4$, the radial and
angular moments give
\[
 \mathbb E[\sigma_2^{-2}\|Q(E)\|_F^2]
 \ge\frac94\varepsilon.
\]
Appending the natural comparator separator $T=\gamma^\dagger$ still
gives
\[
 \mathbb E[\sigma_2^{-2}\|Q(E)T\|_F^2]
 \ge\frac14\varepsilon.
\]
Neither expression is uniformly $O(\varepsilon^2)$.  Yet the full fit
$a\gamma+\varepsilon e\eta$ has range $\operatorname{span}(a,e)$
almost surely, so the actual filtered source, with or without $T$, is
identically zero.  This defeats the proposed proof lemma more strongly
than a badly chosen comparator would, while leaving the leaf learner
successful.  See \sourcefile{claude_screened_residual_audit.md} and
\sourcefile{screened_source_counterexample_audit.md} for the exact
integrals, corrected dimensions, perturbation expansion, and review.

The same audit also rejects two unrelated shortcuts: learned templates
need not be measurable in local probe rows alone, and a per-node bound
by the full residual norm cannot be summed over all tree nodes as though
there were only $L$ of them.  The actual comparator source includes the
signed inherited return as well as the current residual commutator.

Several corrections are needed even before those counterexamples.
The internal predesign is $16k\times64k$, not the retained
$8k\times64k$ design.  Weak directions must be indexed by the smallest
singular values, and $d\le s$ does not imply response rank $d$.
A random inverse singular scale remains inside its joint expectation;
conditioning on a probability-one full-rank event does not make it
constant.  Coefficients formed from $J(B)$ may themselves depend on
exterior $\Psi$ rows.  Finally, with
$B_i=V_i^TJ(B)$, $E_i=V_i^TJ(E)$,
$H=(B_1+E_1)^\dagger$, and $H_B=B_1^\dagger$, the exact expansion is
\[
 Q(E)=E_2-B_2H_BE_1+B_2(H_B-H)E_1-E_2HE_1.
\]
The last two terms cannot be omitted, and the adaptive $V_i$ need not
be comparator-only objects.  These are additional algebraic and
conditioning errors in the relayed proposal, distinct from its false
screened-moment claim.

\section{The old full-range learner: inverse growth and actual failure}
\label{sec:later-fullrange}

Two different obstructions were established for the full-range
commutator candidate with wide-query width $q=64k$, response width
$4k$, retained rank $r=8k$, and physical leaves of size $16k$.  This
candidate retains the entire fitted response and chooses the remaining
directions by inverse-optimal completion.  It is not the maintained
annihilated learner or the later synchronized union learner.  Its
all-depth exact-HSS rank and capture theorem remains valid; the noisy
failure below does not contradict exact recovery.

\subsection{An actual full inverse diverges while approximation is exact}

For $k=1$, let $a,b$ be siblings inside a nonroot node $p$, and take the
fixed one-way rank-one input $A=uv^T$, with $u$ supported in $a$ and $v$
in $b$.  Put $h=\Psi_a^Tu$.  On the $b$ branch every nonzero response
has left factor $h$, despite every outgoing physical block being zero.
For a retained query space $S_j$, full-fit retention implies
\[
 P_{S_j}h=P_{\operatorname{row}(G_j)}h.
\]
Write $\delta_j=\|(I-P_{S_j})\widehat h\|^2$.  Conditional on $h$,
the two phantom children use independent physical Gaussian rows and
have axial symmetry about $\widehat h$.  Their mean residuals equal
$d_\ell\widehat h$, where $d_\ell=\mathbb E\delta_j$.  Comparing the
parent projection with the average of the two child projections gives
\[
 d_{\ell+1}\le\frac{d_\ell+d_\ell^2}{2},\qquad
 d_0=\frac34,\qquad d_\ell\le12\,2^{-\ell}.
\]
The initial value concerns the full sixteen-dimensional Gaussian
predesign: the adaptive eight-dimensional output preserves its fit.

The signal branch represents $u$ exactly.  If $D_a^Tc_a=h$ and
$D_b^Tc_b=P_{S_b}h$, then $\|c_a\|=1$ and
\[
 \sigma_{\min}\!\begin{pmatrix}D_a\\D_b\end{pmatrix}^{\!2}
 \le\|h\|^2\delta_b.
\]
For a range-only tie and canonical-frame rule, scaling $h$ does not
change the selected spaces.  Its Gaussian radius is therefore
independent of $\delta_b$, and
\[
 \mathbb E\|G_p^\dagger\|_{\mathrm{op}}^2
 \ge\frac{1}{(q-2)d_\ell}
 \ge\frac{2^\ell}{12(q-2)}.
 \tag{\thechapter.1}
\]
All the designs in this particular finite-depth family are nonsingular:
the child spaces projected onto $h^\perp$ have independent Haar
orientations and dimensions whose sums are below $q-1$.  This argument
does not declare the selected Grams Wishart.

At $p$, however, both the outgoing block and response are zero.
Inverse-optimal completion may discard the almost-canceling direction;
its retained inverse trace satisfies $V_p\le\min(V_a,V_b)$.  Every
true cut range is captured, and the fixed learned coefficient space
contains $A$ exactly.  Thus (\thechapter.1) refutes a uniform bound for
the \emph{full predesign inverse}, not approximation stability.
The proof, exact scope of the range-only rule, and independent review
are in \sourcefile{phantom_stability.md} and
\sourcefile{phantom_stability_audit.md}.

An associated scalar shortcut also fails.  For a normalized shared
target, define $R=\beta/\|f\|^2$, where $\beta$ is the query residual
and $f$ the minimum-norm physical fit.  At a Gaussian leaf,
\[
 R\ \stackrel{d}{=}
 \frac{X_{q-m}X_{q-m+1}}{X_m},\qquad \mathbb ER=168
 \quad(m=16,q=64),
\]
with independent chi-square variables.  The actual parent law also
depends on child metric matrices and their relative orientations.
Reachable leaf states with an increasingly strong fitted direction and
large optional directions give a limiting conditional mean exceeding
$6/7$ times the child ratio.  Hence the proposed uniform conditional
contraction by $6/7$ is false; this does not rule out the weaker
unconditional estimates subsequently used below.
See \sourcefile{phantom_ratio_route.md} and
\sourcefile{phantom_ratio_audit.md}.

\subsection{A fixed-input lower bound for the entire retained projector}

The later construction controls all seven optional directions, not
only the mandatory fitted line.  For small $\varepsilon>0$, set
\[
 t=\left\lceil\frac72\log(1/\varepsilon)\right\rceil,
 \quad\ell=t+1,\quad N=128\,2^t,\quad L=t+3.
\]
Two physical leaves of size sixteen form a source $a_0$.  Extend it
through $t$ active/blank merges to $a$; let $b$ be its sibling and
$p=a\cup b$ a root child.  With $u=e_1$, $v_b$ a fixed unit vector
uniform on $b$, and $v_{\mathrm{ext}}$ an exterior coordinate, define
\[
 B=\frac{u(v_b+v_{\mathrm{ext}})^T}{\sqrt2},\qquad
 E_0[a_1,a_2]=E_0[a_2,a_1]=I_{16},\qquad
 A_\varepsilon=B+\varepsilon E_0.
\]
All remaining entries of $E_0$ vanish.  The matrix is fixed before
sampling, $B$ is globally rank one, and
$\operatorname{OPT}_1(A_\varepsilon)\le32\varepsilon^2$.
For the full actual retained rank-eight space $U_p$, the reviewed result
is
\[
 \mathbb E\operatorname{dist}(u,U_p)^2
 \ge c\varepsilon^{1.87}.
 \tag{\thechapter.2}
\]
Since $A_\varepsilon[p,p^c]=uv_{\mathrm{ext}}^T/\sqrt2$, every
approximation confined to these learned nested row spaces incurs at
least half this expected squared loss.  Consequently its ratio to
$L\operatorname{OPT}_1(A_\varepsilon)$ diverges.  Refitting or reducing
rank inside the same spaces cannot repair this failure.

The proof has three substantive layers.  First, a fixed compact source
event of probability bounded below gives an initial query gap
$\alpha_0\asymp\varepsilon^2$ and a physical fit within
$O(\varepsilon)$ of the represented signal.  The active/blank physical
martingale has $O(\varepsilon^2)$ squared error uniformly in its length.
Its gap is $\alpha_0$ times $t$ independent
$\operatorname{Beta}(24,4)$ factors.  Since their mean negative
logarithm is less than $4/25$, a high-probability state event satisfies
\[
 c_\alpha\varepsilon^2e^{-(4/25)t}
 \le\alpha\le C_\alpha\varepsilon^2.
\]
For the phantom sibling, the separately proved moment estimates give
on a high-probability event
\[
 \beta\le .439^\ell,\qquad \phi_b\le .699^\ell,\qquad
 V_b\le(9/8)^\ell,\qquad R_b\le(22/25)^\ell.
\]
Here $\phi_b$ is the mandatory inverse cost, $V_b$ the full retained
inverse trace, and $R_b$ the innovation-to-fit ratio.  The fit-decay
bound uses an exactly certified ninth-moment angular estimate, not a
Monte Carlo estimate.  The state event also bounds the total physical
probe Frobenius norm by a constant times $N$.

Second, condition on the complete active state and phantom data
invariant under query rotations fixing $h$.  Eliminate the optional
query coordinates and expose their planes and the remaining radii.
On a further event of conditional probability at least $3/8$, a single
relative angle $c$ remains, with sphere-coordinate density
$g_{49}(c)\propto(1-c^2)^{23}$.  Physical backsubstitution locates a
zero $c_*$ of the \emph{actual physical overlap}; a zero of a projected
query overlap alone would not suffice.  On the interval
\[
 |c-c_*|\le\kappa\frac{\sqrt\alpha}{\sqrt{T_\ell}},
 \qquad T_\ell=C_T(2\cdot .699)^\ell,
\]
the mandatory physical overlap is small and the parent mandatory
inverse cost is at most a constant times $\phi_b$.

Third, the Schur-energy estimate controls every inverse-optimal
padding direction.  Its vanishing terms include
\[
 \varepsilon^2\bigl[(2\cdot .699)(9/8)\bigr]^\ell,
 \qquad [2\cdot .699^2]^\ell,
\]
whose exponents are favorable.  With $\kappa$ fixed sufficiently small,
the full retained space captures at most one quarter of the squared
norm of $u$ throughout the interval.  Its conditional probability is
at least a constant times $\sqrt\alpha/\sqrt{T_\ell}$, hence at least
\[
 c\varepsilon^p,\qquad
 p=1+\frac74\left(\frac4{25}+\log(1.398)\right)
   =1.86632462667\ldots<1.87.
\]
All physical-fidelity and state events are imposed \emph{before} the
final shrinking angle interval.  Subtracting an unrelated
unconditional failure probability from that interval would not prove
the claimed lower bound.

The complete proof and two independent reviews are recorded in
\sourcefile{full_projector_bridge_obstruction.md} and
\sourcefile{full_projector_bridge_audit.md}; its ingredients are
\sourcefile{active_blank_martingale.md},
\sourcefile{phantom_negative_moment_route.md},
\sourcefile{phantom_negative_kernel_bound.md},
\sourcefile{physical_bridge_zero.md}, and
\sourcefile{optional_cost_exclusion.md}.  The theorem concerns exact
full-range selection with the specified completion.  An absolute
singular-value cutoff can change that algorithm at small noise.  The
earlier $\varepsilon^{7/5}$ result concerned only its mandatory line;
it must not be substituted for the full-projector proof.

The maintained annihilated learner survives this particular source
mechanism.  Its phantom sample vanishes and its blank frames are
deterministic.  For the stated leftmost two-leaf source family, with
zero-outgoing siblings and vanishing perturbation commutator above the
source, the reviewed bound is
\[
 \mathbb E\operatorname{dist}_F(A,\mathcal S_\cup)^2
 \le\frac{191}{11}\|E\|_F^2
 \le\frac{764}{33}\operatorname{OPT}_k(A).
\]
These support conditions are essential; this is neither a general
stability theorem nor evidence against (\thechapter.2).
See \sourcefile{maintained_internal_source_bound.md} and
\sourcefile{canonical_blank_path_audit.md}.

\section{Separate views can lose a clean parent signal}
\label{sec:later-separateviews}

Post-training union of independently propagated views does not ensure
that a signal contained in the child union survives the parent union.
The initial oblique example established this only for prescribed local
states.  The later construction reaches it through the actual
maintained and raw two-view rules, including physical leaves, shared
probes, column nullifiers, and canonical completion.

Take $N=32$, $k=1$, retained rank $r=2$, and $q=16$, with physical
leaves of size four.  Let a sixteen-coordinate root child $I$ have
children $a,b$ of size eight.  On the two leaves of $a$ choose
orthonormal coordinate pairs $(u_1,p_1),(u_2,p_2)$ and put
$u=(u_1+u_2)/\sqrt2$.  On $b$ choose coordinates $b_1,\ldots,b_4$;
let $f$ be exterior to $I$.  The fixed input is
\[
 B=uf^T,\qquad E=p_1b_1^T+p_2b_2^T,\qquad
 A_\varepsilon=B+\varepsilon E.
\]
The perturbation is entirely internal to $I$, so $J_I(E)=0$.
For $0<\varepsilon<1$, the singular values of the cut at $a$ are
$1,\varepsilon,\varepsilon$, and
\[
 \operatorname{OPT}_1(A_\varepsilon)=2\varepsilon^2.
\]
Thus the comparator is best, rather than merely a convenient rank-one
matrix.

The explicit full-rank probe witness in
\sourcefile{two_view_union_reachability_route.md} produces child row
planes
\[
 A_a=\operatorname{span}\{(u+p_1)/\sqrt2,p_2\},\quad
 R_a=\operatorname{span}\{(u+p_2)/\sqrt2,p_1\},
\]
and $A_b=\operatorname{span}(b_1,b_2)$,
$R_b=\operatorname{span}(b_3,b_4)$.  Consequently the child union
contains $u$ exactly.  The child sample covariances on
$(u,p_1,p_2)$ can be taken as
\[
 S_{\rm ann}=\begin{pmatrix}2&1&0\\1&2&0\\0&0&3\end{pmatrix},
 \qquad
 S_{\rm raw}=\begin{pmatrix}6&0&2\\0&8&0\\2&0&6\end{pmatrix}.
\]
Their selected projectors have strict retained/discarded gaps.  The
single probe construction realizes both covariances and the required
physical-leaf and column frames simultaneously; they are not prescribed
as independent fictitious responses.

At $I$ both actual row samples have structural rank one.  Their lifted
fits have the forms
$A_1(u+p_1)/\sqrt2+B_1b_1$ and
$A_2(u+p_2)/\sqrt2+B_2b_4$, with all displayed coefficients nonzero.
In each canonical child-template coordinate order, the first $a$
column is therefore independent of the mixed fit.  It is used for
padding before the second $a$ column.  The resulting spaces are
\[
 A_I=\operatorname{span}(u+p_1,b_1),\qquad
 R_I=\operatorname{span}(u+p_2,b_4).
\]
Their union projects $u$ to $(2u+p_1+p_2)/3$, and hence
\[
 \operatorname{dist}(u,A_I+R_I)^2=\frac13.
 \tag{\thechapter.3}
\]
The current residual commutator is zero, the parent outgoing comparator
is rank one, and neither view truncates a nonzero parent sample tail.
The loss is inherited through the separately propagated templates.

This is an actual-recursion reachability statement.  For each fixed
$\varepsilon>0$, full design ranks, spectral gaps, nonzero mixing
coefficients, and the relevant first canonical pivot persist on a
positive-probability Gaussian neighborhood.  Only the needed physical
subspaces and first $a$ pivot must be continuous.  Some unused canonical
columns can change gauge discontinuously, and the proof does not assume
otherwise.  The parent rank-one sample is structural under probe
perturbations, rather than a nongeneric zero singular value imposed by
the witness.

The witness uses narrow-probe coordinates of size $1/\varepsilon$.
Its neighborhood probability may depend on $\varepsilon$; (\thechapter.3)
does not by itself yield an asymptotic expected-error lower bound.
Independent deterministic evaluations of the actual dense learners at
$\varepsilon=1,.1,.01$ reproduce $1/3$ and full rank-four predesigns at
all fifteen nodes.  The algebra and continuity argument, not those
evaluations, prove reachability.  See
\sourcefile{two_view_union_reachability_audit.md} and, for the earlier
unreached local example, \sourcefile{two_view_union_oblique_route.md}.

\paragraph{The synchronized repair and its limits.}
Using the same child-union templates for both parent views removes
this exact mechanism.  The specified synchronized learner retains at
most $4k$ directions with width $32k$, retains the full raw fit on its
low-rank branch, and otherwise combines the lifted raw leading $2k$
directions with the annihilated fit.  Its completion first fills query
visible directions and then physical kernel directions if necessary.
For exact HSS-$k$ and the stated rank-slack inputs, rank and true-cut
containment have proofs.  A clean parent incurs no new loss when its
common child template contains the comparator and its predesign is
injective on the relevant coefficient space; on singular branches an
invisible comparator component must be retained explicitly in the
identity.  Full rank of a retained design is weaker than full rank of
the next stacked predesign.  None of these statements provides the
joint inverse-weighted budget for arbitrary noisy signed sources.
See \sourcefile{synchronized_union_exact_route.md},
\sourcefile{synchronized_union_route.md}, and their audits.

\section{Conditioning, passive memory, and changing targets}
\label{sec:later-conditioning}

\subsection{What the renewed latent law permits one to condition on}

For the independent-source fitted-target tree, the exact innovation
mixture is a theorem for a restricted recursively generated summary.
It may retain physical bases and fits, allowed local and retained
Grams, and deterministic functions of these data.  It may not retain
forgotten query-coordinate designs, old residual vectors or radii, or
arbitrary response cross-Grams.  Selection must be invariant under
right rotations of query coordinates and must contain the entire fit.

For example, after temporarily exposing two current child designs and
their scalar latent variances $\tau_a,\tau_b$, the parent residual
before compression has covariance
\[
 (\tau_a+\tau_b)N_G,
\]
where $N_G$ projects onto the kernel of the full $2r$-row predesign.
The fitted increment and this residual are conditionally independent
Gaussians.  The residual is not isotropic in the larger kernel of the
retained $r$-row design under that same conditioning.  Only after
forgetting the child query frames and using the retained-space
stabilizer can one replace its direction by a uniform direction in
that larger kernel.  The beta--chi-square identity then supplies
\[
 \tau_{\rm new}=(\tau_a+\tau_b)\,Z,\qquad
 Z\sim\operatorname{Beta}\!\left(\frac{q-2r}{2},\frac r2\right)
\]
in the reduced realization.  This is not a license to restore all the
discarded records afterwards.

Two explicit pitfalls illustrate the restriction.  A selector that
inspects a fixed query coordinate can orient an optional direction
toward that coordinate, contradicting the alleged isotropic residual
variance.  Even an invariant record can be excessive: retaining the
old residual radius, the new fit, and its Gram determines the new
radius by subtracting the fitted energy.  A fresh nondegenerate beta
radius is then impossible.  Right-rotation invariance alone is
therefore insufficient.  The corrected theorem and its examples are
in \sourcefile{fitted_target_exact_audit.md}; the matrix version uses
the same restricted-summary obligation in
\sourcefile{matrix_latent_budget.md}.

The source posterior does not renew as a Wishart posterior at every
merge.  With block pre-Gram
$K=\left[\begin{smallmatrix}A&C\\C^T&B\end{smallmatrix}\right]$
and Schur complements $S_a=A-CB^{-1}C^T$,
$S_b=B-C^TA^{-1}C$, the fitted increment covariance is
\[
 \Omega=K^{-1}\operatorname{diag}(\tau_bS_a,\tau_aS_b)K^{-1}.
\]
Its likelihood involves
\[
 \frac{(Kz)_a^TS_a^{-1}(Kz)_a}{\tau_b}
 +\frac{(Kz)_b^TS_b^{-1}(Kz)_b}{\tau_a}.
\]
Already on a scalar cross-block slice this is proportional to
$(\tau_b^{-1}+c^2\tau_a^{-1})/(1-c^2)$, not a linear precision term
in $K$.  Source conjugacy thus does not establish learned-Gram
conjugacy.  Likewise an inverse-weighted residual cost contains the
actual product $\mathbb E[(\tau_a+\tau_b)V_p]$; replacing it by a
product of expectations drops its covariance term.
See \sourcefile{binary_posterior_closure.md}.

These obstacles did not prevent every later local theorem.  The finite
moment ladder bypasses conjugacy, retaining the dependent products and
conditional Haar orientations.  With its stated independent-source
assumptions, $r\ge4096$, $q=8r$, and target width at most $r/2$, it
establishes uniform directional error moments and inverse means with
large absolute constants.  It is not a theorem for $r=2$, nor an
automatic interface to actual exterior HSS commutators.
See \sourcefile{fitted_energy_bootstrap_route.md} and
\sourcefile{fitted_energy_bootstrap_audit.md}.

\subsection{Projection loss, fitted energy, and passive records differ}

A physical projection error can be bounded by the target norm on a bad
event.  A fitted coefficient can be arbitrarily large on that same
event and affect ancestor selections.  Thus paying only the local
projection loss does not dispose of the future fitted-energy cost.
For the fitted-target martingale $S_t=F_t-C_0$, a valid global stopping
argument with exit threshold $M\|C_0\|$ gives
\[
 \mathbb E\|E_H\|^2\le
 \beta+(1+M^{-2})\,
 \mathbb E\|S_{H\wedge\tau}\|^2,
\]
where the stopped quadratic variation includes the exit increment.
Discarding that overshoot changes the process.  A global fit bound does
not normalize every small local target; conditioning on survival can
also couple siblings.  The reduction in
\sourcefile{fitted_target_stopping.md} was valid but did not close its
remaining bound.  Later large-rank source theorems should not be read
back into that scalar stopping argument.

Merely carrying an old target or error vector as a passive record
does not make it an actively retained fit.  The next compression may
discard it, after which its update has an additional projection loss
and is not the centered full-fit recurrence.  Even retaining each
current joint and separate child fit need not preserve older witnesses.
\sourcefile{augmented_memory_route.md} gives a full-rank two-merge
local construction in which current fits and inverse-optimal padding
eventually exclude a nearly represented signal.  That construction
was checked as an admissible local state, not reached from the actual
Gaussian leaf process; it is not an expected-error counterexample.
The joint-plus-split learner nevertheless has a valid exact-input
endpoint at $r=16k$, $q=64k$, response width $4k$, and physical leaves
of size $32k$.  Its three mandatory fits have combined dimension at
most $12k$, and the exact-input boundary-anchor induction proves rank
and cut containment.  A clean rank-one bridge retains a good separate
fit from either child.  That one-bridge protection does not automatically
protect the same vector at the next ancestor.
See \sourcefile{split_fit_learner_audit.md}.

Nor does physical closeness to a fixed envelope suffice.  If
$Q=DC+E$, $D^TE=0$, and $\|E\|_{\rm op}\le\delta<1$, then
\[
 \sigma_{\min}(Q^TX)\ge
 \sqrt{1-\delta^2}\,\sigma_{\min}(D^TX)
   -\delta\|(I-DD^T)X\|_{\rm op}.
\]
This gives good-event inverse moments under a \emph{relative
probe-space} perturbation condition.  Without it, for $n>q$ an adaptive
unit vector arbitrarily close to a fixed $e$ can approach
$\ker X^T$ on a positive-probability event while remaining nonsingular.
Its inverse cost grows like an arbitrary inverse squared parameter.
An accompanying fixed-target example has zero query innovation but
arbitrarily large fitted coefficient.  It can even converge eventually
to the exact core for every realization while its expectations diverge.
These are right-rotation-invariant adaptive constructions, not reached
nested-tree examples.  They refute a bridge based only on geometric
closeness and explain the need for uniform integrability.
See \sourcefile{near_envelope_route.md}.

A related historical capture obstruction is also limited in scope.
The marginal facts $D=T-B$, $T/\sigma^2\sim\chi^2_{15}$, small
$L^2$ residual $B$, and $D\ge_{\rm st}\chi^2_2/2$ alone permit abstract
couplings with an infinite first inverse moment.  Such couplings need
not satisfy the actual fresh latent-radius law.  Once that law is kept,
the later Laplace argument gives
$\mathbb ED^{-p}\le\sigma^{-2p}2^{-p}
\Gamma(3/2-p)/\Gamma(3/2)$ for $p<3/2$, including
$\mathbb ED^{-1}\le\sigma^{-2}$.  Thus the marginal construction is
not an actual-tree counterexample.  Compare
\sourcefile{actual_full_response_capture.md} with
\sourcefile{latent_capture_laplace.md} and its independent audit.

\subsection{Target drops require a retained deterministic guard}

Deterministic compatible transfers
$U_I|_i=U_iA_i$ preserve the additive physical target and transform
omission covariances by congruence,
$C_I=\sum_iA_i^TC_iA_i$.  Invertible changes of output coordinates
preserve the relevant fit ranges.  A strict drop is different: a fixed
physical prefix that would complete the \emph{current} smaller target
may already have been discarded while an earlier larger target was
being fitted.

The explicit two-target source example has a fixed prefix of length
$r-2$ and a random fitted two-plane.  A prescribed next coordinate
almost surely lies outside the retained space.  Dropping one output
at the parent does not recreate that coordinate.  Treating its learned
replacement as a fixed Gaussian direction is invalid.  Without a
history-compatible prefix condition, the inverse decomposition has an
extra nonnegative padding cost
\[
 \mathcal P=\|B_{\rm phys}(I-P_{\operatorname{range}Z})\|_F^2,
\]
which persists when there are more unexcited available directions than
active target columns.  The moment recurrence must then charge
inverse-weighted products involving $\mathcal P$.
See \sourcefile{nested_target_bootstrap_route.md}.

The fixed-guard modification repairs this particular problem.  With
$s_{\max}\le r/2$, put $c=r-s_{\max}$.  Each source has a deterministic
pool of $2c$ directions and retains its first $c$ as a guard.  At a
parent, the block sum of its two child guards is an available
deterministic pool of $2c\ge r$ directions; select the full fit and
complete by scanning this pool.  Since the fit rank is at most
$s_{\max}$, the first $c$ pool directions remain retained.  This is an
implementable induction, not an oracle reopening a discarded physical
coordinate.  The enlarged-core argument then restores the original
local moment constants under arbitrary compatible deterministic drops.
See \sourcefile{fixed_guard_completion_route.md} and
\sourcefile{fixed_guard_completion_audit.md}.  It repairs a local source
tree; independent innovations for arbitrary HSS assembly remain a
separate requirement.

\section{Exterior information and pilot independence}
\label{sec:later-exterior}

\subsection{Comparator anchors do not describe actual residual records}

At a fixed cut $I$, all exterior \emph{comparator} anchors restrict
inside $I$ to its rank-$k$ outgoing space.  Conditioning on those
anchors removes at most $k$ local Gaussian directions.  Descendant
anchors are not included in this statement.  An orthogonal comparator
wavelet decomposition also does not make children conditionally
independent: if
\[
 \begin{pmatrix}X_a\\X_b\end{pmatrix}\Bigm|X_I
  =\begin{pmatrix}R_a\\R_b\end{pmatrix}X_I+QG,
\]
then $\operatorname{Cov}(X_a,X_b\mid X_I)=-R_aR_b^T$ in normalized
coordinates.  Fixing a parent sum creates a contrast coupling.

Actual exterior sketches can expose arbitrarily more directions, even
with a best comparator.  At a root cut $K,J$, take fixed rank-$k$
frames $U,V$ and orthonormal leaf-supported directions
$P=(p_1,\ldots,p_M)$, $Q=(q_1,\ldots,q_M)$ orthogonal to the respective
comparator frames.  Set
\[
 B=\sigma(UV^T+VU^T),\qquad E=\varepsilon P Q^T,
 \qquad 0<\varepsilon<\sigma.
\]
On the specified physical-leaf/coarse tree, $B$ is HSS-$k$, and the
root-cut singular values give
$\operatorname{OPT}_k(B+E)=M\varepsilon^2$ exactly.  Fix the narrow
probe, exterior wide rows, and comparator anchor.  Each actual exterior
leaf commutator then reveals
$\varepsilon h_jx_j$, where $h_j=\Psi_K^Tp_j$ and the observed
row $x_j=q_j^T\Omega_{J_j}$ is nonzero almost surely.  It therefore
reveals every $h_j$.  The remaining local physical Gaussian covariance
is $I-UU^T-PP^T$, losing $M$ residual directions, not just $O(k)$.
This is a conditioning obstruction, not an approximation lower bound.
See \sourcefile{correlated_source_interface_route.md} and its audit.

The same distinction applies to internal injections.  Deterministic
additions in the currently retained guard pool are reproduced exactly
and preserve the local theorem.  A genuinely omitted addition is fresh
only if its physical query directions were reserved from the entire
previous computation, including target responses and conditioning
records.  Orthogonality to the final selected template alone does not
give that independence.  Ordinary binary parent rows already belong to
the children; comparator contrasts of dimension at most $2k$ do not
automatically supply unused Gaussian coordinates.  Even a valid
reserved-source model must charge injections by their ages, since
recent injections have not undergone many residual contractions.
See \sourcefile{guarded_internal_injection_route.md}.

\subsection{Independent pilots leave exterior dependence and width costs}

Freeze an independently learned rank-$r$ column pilot $V_I$, then use
fresh probes and its annihilator $N$.  With
$R=A[I^c,I](I-V_IV_I^T)$ and annihilator width $s$,
\[
 Z=\Psi_{I^c}^TR\Omega_I N,\qquad
 \mathbb E(\|Z\|_F^2\mid V_I)=qs\|R\|_F^2.
\]
For a fixed $m$-dimensional local template the corresponding lifted
mean is $ms\|R\|_F^2/(q-m-1)$.  These are useful exact one-step costs.
But a fixed incoming cut of rank $r+1$ defeats every rank-$r$ pilot.
Conditioning on that pilot and the local design, a nonzero surviving
incoming response produces a Gaussian fit with full physical covariance.
A retained space containing it cannot be measurable in local fresh
rows alone.  Exposing a common exterior anchor subsequently correlates
siblings.  Independence of the pilot from the fresh probes does not
prove local independence of the \emph{new learned} hierarchy.
See \sourcefile{crossfit_hss_interface_route.md} and its audit.

An oracle collection of independent full-row exterior pilot ranges is
also not automatically a bounded-width nested envelope.  Along a path,
fixed shell perturbations can add $p$ mutually orthogonal directions
at each of $h$ levels.  Differences of ancestor sketches expose each
new block through an invertible $p\times p$ Gaussian factor, forcing
dimension at least $hp$ in the nested closure at the terminal node.
Conversely, intersecting a clean parent signal envelope with a noisy
child range can yield the zero space.  Projecting that signal into the
child range restores feasibility but reintroduces adaptive dependence.
These facts in \sourcefile{pilot_route.md} do not exclude all pilots;
they exclude those proposed exact-envelope constructions.

\subsection{Clean full-row linear extraction costs depth}

Let $W\in\mathbb R^{N\times M}$ collect every forward query.  An
input-independent linear extraction of a clean exterior sketch at $I$
has the form $A[I,:]WK_I$ with $W[I,:]K_I=0$.  On a path
$I_0\supset\cdots\supset I_L$, choose a fixed rank-$k$ test block
between $I_d$ and $k$ coordinates of its sibling.  Exact capture of
that test forces
\[
 \operatorname{rank}W[I_{d-1},:]
  -\operatorname{rank}W[I_d,:]\ge k.
\]
Summation gives $M\ge kL$.  Only finitely many fixed test inputs are
needed for the simultaneous probability-one assertion.  Adding
transpose queries of width $M_H$ does not evade it on nodes larger
than $M_H$: an invisible perturbation
$(I-P_H)T(I-P_W)$ shows that exact extraction either has its right
factor in $\operatorname{range}W$, or requires the entire coordinate
space of $I$ in $\operatorname{range}H$.

Finite sign averaging at one partition has the additional problem
that its variance is comparator-sized, even when the approximation
residual is zero; finite Rademacher masks can erase a fixed rank-one
interaction with positive probability.  The lower bound concerns clean
\emph{full-row}, input-independent linear extraction.  The two-sided
compressed commutator already available with $O(k)$ probes lies outside
that model, as do nonlinear and adaptive decoders.
See \sourcefile{tree_mask_route.md}.

\section{Frozen compression and shared external noise}
\label{sec:later-frozen}

For independent pilot frames $P,Q$ and fresh Gaussian probes, frozen
compression has the exact conditional law
\[
 Y_c=A_c\Omega_c+N_y,\qquad X_c=A_c^T\Psi_c+N_x,
 \quad A_c=P^TAQ,
\]
where the four arrays $\Psi_c,\Omega_c,N_y,N_x$ are initially
independent.  The noise columns have fixed, possibly dense physical
covariances $C_y,C_x$, and
\[
 b=\operatorname{tr}C_y+\operatorname{tr}C_x
 \le\|A-PP^TAQQ^T\|_F^2.
\]
If the pilot contains a comparator, $b\le\|E\|_F^2$.  At each complete
level the expected commutator-noise energy is exactly
$q_\psi q_\omega b$.  These are orthogonal energy identities, not
independence between nodes; reusing the probes in learned subsequent
compressions does not renew the four-array law.
See \sourcefile{frozen_compression_noise_route.md} and its audit.

Factor $C_y=L_yL_y^T$, $C_x=L_xL_x^T$.  Every actual core computation
can be realized simultaneously as an exact commutator computation for
\[
 A_{\rm ext}=\begin{pmatrix}A_c&L_y\\L_x^T&0\end{pmatrix},
 \quad\Psi_{\rm ext}=\begin{pmatrix}\Psi_c\\G_x\end{pmatrix},
 \quad\Omega_{\rm ext}=\begin{pmatrix}\Omega_c\\G_y\end{pmatrix},
\]
with independent standard Gaussian exterior arrays.  The external
perturbation has squared Frobenius norm $b$.  This pathwise realization
preserves the actual core recursion; its root is a proper node of the
augmented problem, not a new zero-commutator global root.

Rank-one incoming noise already defeats the independent-source product
law.  Set $A_c=N_y=0$ and $N_x=vh^T$.  The physical-leaf row frames
are deterministic, so two first internal sources have independent
Gaussian $M$-row designs $G_a,G_b$ but the same $h$.  Their residuals
$e_i=(I-P_{G_i})h$ obey
\[
 \begin{aligned}
 B_i&=\|e_i\|^2=\|h\|^2Z_i,\\
 Z_a,Z_b&\ \hbox{independent }
     \operatorname{Beta}((q-M)/2,M/2),\\
 \operatorname{Cov}(B_a,B_b)&=\frac{2(q-M)^2}{q}>0.
 \end{aligned}
\]
Each marginal is the expected chi-square law.  At $q=8r,M=2r$ the
covariance is $9r$.  Exposing the complete source designs and fits gives
$h=\mu+\eta$, where, conditionally on the designs,
$\eta\sim N(0,I-P_{[G_a;G_b]})$ is independent of the fits.  The next
retained design $R$ lies in the old stacked row space, so
\[
 \|(I-P_R)h\|^2
 =\|(I-P_R)\mu\|^2+\chi^2_{q-2M}
\]
under this complete record.  Ignoring the coherent mean is incorrect.
This stronger conditioning is not a contradiction to the restricted
latent theorem; the independent-source hypothesis already fails by the
positive unconditional covariance.  Nor does the example prove physical
instability, since its core matrix is zero.
See \sourcefile{independent_noise_tree_route.md} and its audit.

There is a positive rank-slack exception.  If augmented outgoing and
incoming cut ranks fit within the maintained retained rank, its own
hybrid column hierarchy captures the relevant external incoming factors,
and its nullifiers remove them.  In particular a rank-one incoming
external perturbation of an exact HSS-$k$ core can leave the maintained
row trajectory identical to the noiseless trajectory.  The hybrid's
own column frames are essential; substituting raw companion frames
does not prove the assertion.  Low covariance trace alone does not
imply these rank inequalities.  See
\sourcefile{two_view_external_noise_route.md} and its audit.

The noncentral projector theorem supplies another one-step estimate,
without pretending inherited noise is fresh.  For fixed $M$, standard
Gaussian $Z\in\mathbb R^{t\times s}$, $s<t$, and $U=\operatorname{range}M$,
\[
 \mathbb E P_{M+Z}\preceq P_U+\frac st P_{U^\perp}.
\]
The coefficient $s/t$ is sharp; singular covariance is handled by
quotienting out deterministic mean-only columns before whitening.
There is no general opposite inequality
$\mathbb E P_{M+Z}|_U\succeq(s/t)I_U$: for
$t=3,s=2,M=[a e_1,\delta e_2]$, first large $a$ then small positive
$\delta$ gives limiting leverage $1/2$ on $e_2$, below $2/3$.
In a guarded injection this upper bound interpolates an inherited-fit
metric cost with a full pretemplate cost.  It does not bound the
inherited cost or close arbitrary mixed histories.
See \sourcefile{noncentral_projector_route.md} and its audit.

\section{Consistency correction and global correction barriers}
\label{sec:later-compression}

\subsection{Consistent observations can require large matrix noise}

Noisy compressed products generally violate
$\Psi^TY=X^T\Omega$.  Orthogonal least-squares repair is explicit.
With $S=\Psi\Psi^T$, $T=\Omega\Omega^T$, its minimum-norm matrix
representative solves
\[
 SM_*+M_*T=\Psi X^T+Y\Omega^T.
\]
If the noise-free core is $A$, then
\[
 M_*=A-(I-P_\Psi)A(I-P_\Omega)+Z,
 \quad SZ+ZT=\Psi N_x^T+N_y\Omega^T.
\]
The repaired observations are products of $A+Z$.  The accurate witness
$A+Z$ need not be computable from underdetermined observations; the
computable $M_*$ loses the invisible block.  These two matrix objects
must not be identified.

In eigenbases of $S,T$, conditional noise energy equals
\[
 \mathbb E(\|Z\|_F^2\mid\Psi,\Omega)
 =\sum_{\lambda_i+\mu_j>0}
   \frac{\lambda_i c_{x,j}+\mu_j c_{y,i}}
        {(\lambda_i+\mu_j)^2}.
\]
Away from square dimensions, inverse-Wishart moments give a bound by
$b$ when each larger dimension is at least twice the smaller plus one.
Near-square behavior is different.  With $n_r=2q$, $n_c=q+2$, both
query widths $q$, and $C_x=0$, the mean correction energy is at least
$qb/2$.  A physical nullspace on the opposite side and dimension gap
at most one gives an infinite mean.  In the doubly square case the
mean is also infinite: the smallest two eigenvalues are bounded above
by independent $\chi^2_1$ variables, whose sum has an infinite inverse
mean.  This defeats an unqualified stable orthogonal consistency
repair, not every regularized estimator.  Even where finite, $Z$ depends
on the same designs and cannot be fixed while those designs are
subsequently treated as fresh.
See \sourcefile{compression_consistency_repair_route.md} and its audit.

\subsection{A few global low-rank stages need a global-tail premise}

For a known pilot $M_0$, fresh range queries on $A-M_0$ and adapted
transpose queries give an implementable correction
\[
 M_1=M_0+QQ^T(A-M_0),\qquad
 \mathbb E(\|A-M_1\|_F^2\mid M_0)
 \le\left(1+\frac{k}{p-1}\right)\tau_k(A-M_0),
\]
where $Q$ comes from a Gaussian sample of width $k+p$, $p\ge2$, and
$\tau_k$ is the squared global rank-$k$ tail.  For $p=k+1$ the total
new query cost is at most $4k+2$.  The missing premise is control of
that \emph{global} tail by $L\operatorname{OPT}_k$.

HSS rank does not supply it: $m$ disjoint unit rank-one sibling blocks
are HSS rank one but have global tail $m-k$.  A fixed number of
rank-$O(k)$ corrections from zero cannot recover them when $m$ is
large.  Similarly, with frozen $P,Q$,
\[
 \|A-PCQ^T\|_F^2
 =\|A-PP^TAQQ^T\|_F^2+\|P^TAQ-C\|_F^2;
\]
exactly solving the compressed core cannot improve the fixed geometric
floor.  Requerying full products at each level costs the rank of the
concatenated lifted query spaces, which can be proportional to the
number of levels.  Preserving free leaf-diagonal blocks is also
necessary: an identity matrix is HSS rank zero but has large global
rank.

Even orthogonal projection of an HSS-$k$ comparator onto an arbitrary
fixed nested coefficient space need not retain HSS rank $k$: different
levelwise projections can make formerly common row factors independent
on one cut.  Exact coefficient projection is not automatically
available from small fixed transcripts either.  An invisible-matrix
test shows that exact projection onto a space containing every
diagonal matrix requires
$\dim\operatorname{range}\Psi+\dim\operatorname{range}\Omega\ge N$.
These statements do not invalidate the separately proved approximate
independent refit, nor an offline best-HSS rank reduction once a good
matrix approximation has actually been obtained.  They distinguish
query complexity from a fast or directly computable optimization rule.
See \sourcefile{bounded_stage_compression_route.md} and its audit.

\paragraph{Unreviewed draft extension at the research stop.}
The final \sourcefile{replicated_global_tail_obstruction.md} was saved
as a proved draft but had not received an independent audit.  It
should not be assigned the review status of the preceding results.
Its pointwise compression argument claims that, for disjoint tested
cuts with losses $e_j$ and amplitudes $a_j$, every pilot obeying their
learned row constraints satisfies
\[
 \tau_d(A-M_0)\ge
 \sum_j|a_j|^2e_j-\sum_{\text{largest }d}|a_j|^2e_j.
\]
It proposes applying this to block-diagonal copies of the reviewed old
full-range counterexample, obtaining a failure of any fixed number of
global correction directions to restore an $O(L\operatorname{OPT})$
tail premise for that learner.  It does not target the synchronized
learner or arbitrary nonlinear structured corrections.  No new claim
of independent verification is made here.

\section{Numerical evidence and the final status map}
\label{sec:later-status}

The saved experiments test different objects and should not be pooled
into a single claim.  The old full-range proof-chain batch contained
1064 recursive cases up to $N=65536$, with 32--160 independent draws
per noisy cell.  It identified the bridge mechanism and checked
identities, but the later analytic rare-angle argument proves the
$\varepsilon^{1.87}$ lower bound.  A conditional rotation diagnostic
on one child state is likewise not an unconditional expectation bound.
The guarded full-range batch contained 800 cases from 160 independent
probe pairs, reusing pairs across noise amplitudes.  At $N=131072$ and
$\varepsilon=.1$, mean signal loss rose from $.08175$ to
$.62016\pm.05154$ across the bridge.  This finite amplification neither
refutes an $O(L)$ estimate nor tests the large-rank independent-source
theorem's hypotheses.  See \sourcefile{proof_chain_numerical_summary.md},
\sourcefile{bridge_conditional_rotation.md}, and
\sourcefile{guard_fullrange_bridge_check.md}.

The synchronized learner's separate batch used 624 cases from eighty
independent probe pairs.  Its largest exact squared coefficient-space
error was $3.162\times10^{-23}$; its largest noisy-cell mean error
divided by the \emph{specified comparator residual} was
$1.00792492\pm.00335957$.  The high-rank internal union branch was
exercised 10752 times.  The experiment used an oracle projection onto
the fixed learned coefficient space, not the final random refit or a
rank-$k$ output, and did not compute best-HSS OPT.  Reused seeds across
families and amplitudes are paired cases, not extra independent draws.
The follow-up signed diagnostic reused three probe pairs and checked
400 reached oriented nodes.  Its maximum vector-identity discrepancy
was $2.721\times10^{-15}$, with both substantial cancellation and
amplification between the source and inherited-return terms.  No
expectation factorization follows from those checks.
See \sourcefile{synchronized_union_numerical_check.md},
\sourcefile{synchronized_union_numerical_audit.md}, and
\sourcefile{synchronized_signed_identity_check.md}.

The noncentral projector experiment is narrower still: fixed means,
covariances, and positive semidefinite metrics with fresh Gaussian
arrays.  Its 1150000 evaluations include 20000 repeated noiseless
deterministic evaluations, so only 1130000 are nonempty Gaussian draws.
Scalar and central equality checks, singular quotient identities, and
the anisotropic negative control agree with the proved formulas within
their recorded numerical and sampling errors.  It does not test
freshness of inherited HSS noise or closure of the mixed-history cost.
See \sourcefile{noncentral_projector_numerical_check.md} and its audit.

The endpoint can be stated without conflating these scopes:
\begin{enumerate}[leftmargin=*,label=\arabic*.]
 \item \textbf{Defeated selectors and acquisition models.}
 Minimum-norm and the stated nuclear/tree penalties fail even with
 exact leaf blocks.  The one-sided Gaussian pilot and clean full-row
 linear extraction have their stated lower bounds.  None is an
 impossibility theorem for all nonlinear two-sided decoders.
 \item \textbf{A defeated actual learner.}
 The old inverse-optimal full-range candidate fails the desired
 expected physical bound.  Its earlier mandatory-line padding gap was
 closed by the full-projector proof.  The maintained learner's positive
 theorem on that source family remains compatible with this result.
 \item \textbf{A reached local defect repaired by synchronization.}
 Separate-view propagation can incur a clean-parent loss despite zero
 child-union loss.  A common-template synchronized merge removes that
 specific defect under its visibility hypothesis.  Its exact and
 rank-slack endpoints do not prove noisy physical stability.
 \item \textbf{Local source issues with precise repairs.}
 Restricted conditioning restores the latent law; fixed guards repair
 historical target-drop padding; the finite moment ladder closes the
 sufficiently large-rank independent-source model.  These repairs do
 not manufacture independent external commutator sources.
 \item \textbf{Remaining coupled quantities.}
 Shared external noise, residual-dependent anchors, and the adaptive
 left lift remain coupled to inherited physical errors.  The proved
 unlifted signed-source and tail budgets must still be converted into
 a global expected physical budget without multiplying by an
 independently averaged inverse or double-counting cumulative losses.
 \item \textbf{Reconstruction and the unresolved conjecture.}
 Independent fixed-space refitting and query-only rank reduction remain
 available under their proved premises.  No documented result supplies
 the required all-input synchronized basis-quality premise.  Numerical
 checks do not supply it, and the final replicated-tail note remains
 an explicitly unreviewed draft extension.  The original general
 $O(k)$-query, $O(L\operatorname{OPT}_k)$ conjecture is unresolved.
\end{enumerate}

\appendix
\chapter{Source provenance}
The following local records supply the result statements, their assumptions,
and the internal audits cited in this manuscript. They are provenance for
the present research record, not public bibliographic publications.
The assembled LaTeX file is self-contained. Filenames below identify the
original supporting records in the project directory.

\begin{itemize}[leftmargin=1.3em,itemsep=1pt,topsep=3pt]
\item \sourcefile{active_blank_martingale.md}
\item \sourcefile{actual_full_response_capture.md}
\item \sourcefile{augmented_memory_route.md}
\item \sourcefile{binary_posterior_closure.md}
\item \sourcefile{bounded_stage_compression_route.md}
\item \sourcefile{bridge_conditional_rotation.md}
\item \sourcefile{canonical_blank_path_audit.md}
\item \sourcefile{certified_selector_audit.md}
\item \sourcefile{certified_selector_route.md}
\item \sourcefile{claude_screened_residual_audit.md}
\item \sourcefile{compression_consistency_repair_route.md}
\item \sourcefile{correlated_source_interface_route.md}
\item \sourcefile{crossfit_hss_interface_route.md}
\item \sourcefile{fitted_energy_bootstrap_audit.md}
\item \sourcefile{fitted_energy_bootstrap_route.md}
\item \sourcefile{fitted_target_exact_audit.md}
\item \sourcefile{fitted_target_stopping.md}
\item \sourcefile{fixed_guard_completion_audit.md}
\item \sourcefile{fixed_guard_completion_route.md}
\item \sourcefile{frozen_compression_noise_route.md}
\item \sourcefile{full_projector_bridge_audit.md}
\item \sourcefile{full_projector_bridge_obstruction.md}
\item \sourcefile{guard_fullrange_bridge_check.md}
\item \sourcefile{guarded_internal_injection_route.md}
\item \sourcefile{hss_audit_report.md}
\item \sourcefile{hss_complete_proof_chain.tex}
\item \sourcefile{hss_global_selector_obstruction.tex}
\item \sourcefile{independent_noise_tree_route.md}
\item \sourcefile{latent_capture_laplace.md}
\item \sourcefile{maintained_internal_source_bound.md}
\item \sourcefile{matrix_latent_budget.md}
\item \sourcefile{near_envelope_route.md}
\item \sourcefile{nested_target_bootstrap_route.md}
\item \sourcefile{noncentral_projector_numerical_check.md}
\item \sourcefile{noncentral_projector_route.md}
\item \sourcefile{one_sided_gaussian_pilot_audit.md}
\item \sourcefile{one_sided_gaussian_pilot_obstruction.md}
\item \sourcefile{optional_cost_exclusion.md}
\item \sourcefile{phantom_negative_kernel_bound.md}
\item \sourcefile{phantom_negative_moment_route.md}
\item \sourcefile{phantom_ratio_audit.md}
\item \sourcefile{phantom_ratio_route.md}
\item \sourcefile{phantom_stability.md}
\item \sourcefile{phantom_stability_audit.md}
\item \sourcefile{physical_bridge_zero.md}
\item \sourcefile{pilot_route.md}
\item \sourcefile{proof_chain_numerical_summary.md}
\item \sourcefile{replicated_global_tail_obstruction.md}
\item \sourcefile{screened_source_counterexample_audit.md}
\item \sourcefile{selector_with_leaf_data.md}
\item \sourcefile{split_fit_learner_audit.md}
\item \sourcefile{synchronized_signed_identity_check.md}
\item \sourcefile{synchronized_union_exact_route.md}
\item \sourcefile{synchronized_union_numerical_audit.md}
\item \sourcefile{synchronized_union_numerical_check.md}
\item \sourcefile{synchronized_union_route.md}
\item \sourcefile{tree_mask_route.md}
\item \sourcefile{two_view_external_noise_route.md}
\item \sourcefile{two_view_union_oblique_route.md}
\item \sourcefile{two_view_union_reachability_audit.md}
\item \sourcefile{two_view_union_reachability_route.md}
\end{itemize}
\begin{thebibliography}{9}
\addcontentsline{toc}{chapter}{Public references}
\bibitem{amsel2025}
N. Amsel, T. Chen, F. Duman Keles, D. Halikias, C. Musco, C. Musco, and D. Persson.
\emph{Quasi-optimal hierarchically semi-separable matrix approximation}.
arXiv:2505.16937v2, 2025.
\url{https://arxiv.org/abs/2505.16937v2}.
\bibitem{levitt2024}
J. Levitt and P.-G. Martinsson.
\emph{Linear-complexity black-box randomized compression of rank-structured matrices}.
SIAM Journal on Scientific Computing, 2024; arXiv:2205.02990v3.
\url{https://arxiv.org/abs/2205.02990v3}.
\end{thebibliography}

\end{document}
